% =====================================================================
%  R. Nielsen, OLD AND NEW PROPHETS.  Six Lectures.  Copenhagen: Gyldendal (F. Hegel & Søn), 1878.  112 pp.
%  (Gamle og nye Propheter.  Sex Forelæsninger.)
%  ENGLISH TRANSLATION, from transcription.tex (Danish).
%  GENERATED by ebook-method/gamle-nye-profeter/tr_build.py --write from this template and the six translators' parts
%  in ebook-method/gamle-nye-profeter/tr/ (TRANSLATOR-BRIEF.md; terms: GLOSSARY.md, the title and four key terms
%  approved by Hans 2026-10-12).  Conventions: TRANSLATION-PLAYBOOK.md.  Page numbers of the 1878 print in the margin.
% =====================================================================
\documentclass[12pt,a4paper,openany]{book}
\usepackage[utf8]{inputenc}
\usepackage[T1]{fontenc}
\usepackage{libertinus}
\usepackage{libertinust1math}
\usepackage[english]{babel}
\usepackage[protrusion=true,expansion=true]{microtype}
\usepackage{setspace}
\usepackage[margin=1.3in]{geometry}
\usepackage{fancyhdr}
\usepackage{hyperref}
\hypersetup{pdftitle={Old and New Prophets}, pdfauthor={Rasmus Nielsen}, hidelinks}
\setlength{\headheight}{15pt}
\pagestyle{fancy}
\fancyhf{}
\fancyhead[LE,RO]{\thepage}
\fancyhead[C]{\textit{R. Nielsen: Old and New Prophets (1878)}}
\setstretch{1.3}
\usepackage{marginnote}
\renewcommand{\thefootnote}{*)}
\newcommand{\opage}[1]{\marginnote{\footnotesize\textit{[#1]}}}
% a lecture opens a new page: heading, a short rule
\newcommand{\lecture}[1]{\clearpage\begin{center}{\large #1}\\[0.4em]\rule{3em}{0.4pt}\end{center}\medskip\noindent}

\begin{document}

% PDF 8, the title page
\begin{titlepage}
\begin{center}
\vspace*{3em}
{\Huge OLD AND NEW PROPHETS.}\\[2.5em]
SIX LECTURES\\[1em]
{\small BY}\\[1em]
{\large R. NIELSEN.}\\[5em]
COPENHAGEN.\\[0.5em]
GYLDENDAL'S BOOKSHOP (F. HEGEL \& SØN).\\[0.5em]
{\small PRINTED BY J. JØRGENSEN \& CO.}\\[0.5em]
{\small 1878.}\\[3em]
{\small Translated from the Danish}
\end{center}
\end{titlepage}

% PDF 10, the contents (no page numbers in the print)
\begin{center}{\large Contents.}\\[0.3em]\rule{3em}{0.4pt}\end{center}
\noindent\emph{First Lecture:} The Calling and Position of the Prophets in Past\par
\noindent\hspace*{2em}and Present.\par
\begin{center}\textbf{Old Prophets.}\end{center}
\noindent\emph{Second Lecture:} Moses, Samuel, Elijah.\par
\noindent\emph{Third Lecture:} Prophetic Foresight; Images of the Messiah.\par
\noindent\emph{Fourth Lecture:} The Fulfilment of the Prophecies; Conflicts\par
\noindent\hspace*{2em}between the Old and the New.\par
\begin{center}\textbf{New Prophets.}\end{center}
\noindent\emph{Fifth Lecture:} Saint-Simon, Fourier, Proudhon: Denial of\par
\noindent\hspace*{2em}God and Providence; Nihilism.\par
\noindent\emph{Sixth Lecture:} Reckoning; Neo-Rationalism.\enspace Conclusion.\par
\clearpage

% --- p. 1 ---
\setcounter{footnote}{0}%
\lecture{First Lecture.}\opage{1}As the great teachers of mankind, as the leaders of the peoples
forward toward the promised land of hope, as the torch-bearers of the spirit
on dark, untrodden paths, the prophets hold a high rank. But --- one asks --- are these
leaders now also reliable? The old prophets are, after all,
dead, and the visions of the new ones come at times, here and there,
in disquieting confusion. Is it really a law of history
that every spiritually stirred age must have
its prophets? A critical comparison between antiquity
and the present might at the outset tempt one to suppose
the opposite.

The ancient world, poor in experience but rich
in presentiments and inspired promptings, groped in an outward
half-darkness. For want of secure laws and reliable
facts, the inner lights, with
visions and revelations, had to be its foremost guiding stars:
the ancient world needed prophets. But
the new world, with its armed eye, its telescopes
and microscopes, proceeds by the way of calculation and
requires another guidance. In politics, in political economy,
in industry and trade one does not consult
oracles; in a financial crisis one seeks advice not from
prophets, but from expert professionals.

% --- p. 2 ---
\setcounter{footnote}{0}%
\opage{2}The prophet is like a seer on the mountain. Far off on the
horizon, images of the future emerge from the mist of possibilities
and are drawn in floating outline upon the airy
ground of the clouds; it is refractions of the inner
light that give these figures life and color and full
presence. Here there is no reliable measure of distance;
the prophet sees what is to come as something present:
he does not calculate, he beholds. The reasoning inquirer,
on the other hand, arrives at his cognitions piecemeal. For
him what is to come is a consequence, determined by law according to the connection of things,
of what is past; he does not rely
on a merely intuiting imagination; he doubts and
tests, he reflects and calculates. The latter seems
to the enlightened of the present to be fully as sure as
the former. True enough, we cannot calculate in advance
a political upheaval as the astronomers calculate an
eclipse of the sun; unforeseeable circumstances will always
play a part in historical events; the most carefully
calculated future usually has some surprise or other
in reserve. But even if we cannot
predict exactly what must necessarily happen, we
can nonetheless find out by calculation what will in all
probability happen. The calculus of probability, which
in our day has gained a wide domain, shows its
superiority in the most diverse fields. If
one has a sufficient number of statistical data on a country's
population over a longer series of years, then
one can, without prophetic inspiration, calculate with fair
certainty how many deaths will occur
in the coming year, how many suicides, how many
will hang themselves, how many will drown
themselves, shoot themselves, etc. As regards marriages,
% --- p. 3 ---
\setcounter{footnote}{0}%
\opage{3}one calculates not merely how many will marry in the course of the year,
but even how many with young girls, how many
with older girls and how many with widows. So far
have things come that the postal service can calculate how many
letters will be handed in within a year and a day
with a wrong address and how many without an address,
because in one's haste one has forgotten the address. The calculus of probability
is in steady progress; the statistical
tables grow in many-sidedness as well as in scope, and
will, as far as the educated are concerned, soon, it may be, be said
to have more or less replaced the prophecies.

This consideration is misleading, for it displaces
the point of view.

In the outward world, indeed even in the inward, there is
much that can be calculated, and anyone who in such
things will guess his way forward, or rely on presentiments, or
choose blindly, must take the consequences of having, out of indolence,
neglected to think and calculate. But with
the highest concerns of the spirit it is something else.
Should there come a time in which the prophetic
had no significance, then it would surely have to be a time
in which the poetic, in the deeper sense, had no significance
either, a time where imagination and feeling, the stirrings
of the heart and ideal longings, had no significance,
a time in which the soul, the conscience, with everything personal,
spent its powers on the acquisition of finite
things and the calculation of utility and advantage. In such a
time religion would certainly be a fancy,
faith superstition, the majesty of the spirit dethroned. But
such a time has not yet come and will, it is to be hoped,
never come.

% --- p. 4 ---
\setcounter{footnote}{0}%
\opage{4}It is being prepared, though, people say. Long enough has
the old oriental fable of the arbitrary power of a miracle-working
spirit tyrannized human
nature; it is time that this nature be allowed to
move freely and create for itself its own forms of existence.
The spirit is not a being by itself, exalted
above nature; it is nature itself which in man
comes to consciousness and becomes spirit. All human
relations, private as well as public, are determined
by natural demands and have their ultimate ground in inescapable
natural conditions. The man's love for the woman,
the mother's love for the child, the people's love for
the fatherland, what is that but sheer natural determinations?
The same holds of beauty and
honor, indeed of truth and right. He who does not have an
original, i.e.\ a natural, sense of the beautiful
never comes to know what beauty is. If a
human being has by nature no sense of honor, what use is
it then to try to inculcate in him the principles of honor?
By suitable application of punishment and reward
one can perhaps force unruliness into the forms of a prescribed
morality; one can, after all, break in the refractory
individual as one breaks in the wild horse, by letting
him thirst. But does that promote morality itself? Even if
one could, with the enormous apparatus of the wrath and grace of a supernatural
arbitrariness, hold out the prospect of nine hells
on the one side and as many heavens on
the other, one would still not be able, by outward
pressure, to press true recognition of truth,
freedom and right into a single human soul.

There is an indissoluble bond between nature and
truth. To deny what follows from the
% --- p. 5 ---
\setcounter{footnote}{0}%
\opage{5}nature of a matter is to deny the truth; to keep strictly
to the nature of the matter is to keep to the truth. A great
fuss is made about sin; people weep solemnly over
fallen nature. But what is it that produces
sin and corrupts nature? It is arbitrariness.
Nature in all naturalness cannot corrupt itself;
it is precisely the arbitrary, which emerges as natural life
is to pass over through man into the life of the spirit,
that produces sin. As arbitrariness misjudges
the demands of nature and works against its striving,
the false, the unnatural, comes forth. Arbitrariness
is the principle of evil, the real cause of
discord and disorder. What is disorder but
unnature? By arbitrary precepts one would force
persons of unlike minds, of conflicting natures, to go on
living together in marriage because they now once
are in marriage, to live as if they were at one at heart,
because arbitrariness has commanded that they shall be
at one, to live so that it looks as though they loved
each other inwardly, although they are manifestly close to hating
each other to the death; this untruth, this soul-consuming
unnature shall --- so runs the command of arbitrariness --- be
the right sign of moral order, the genuine proof of
the sanctity of marriage. What is it that brings
society into disorder? This unnatural thing, that something which
by its essence is of a lower kind is set up as a
higher, while what is by nature higher is held under a
pressure that stamps it as a lower; this unnatural thing,
that many must suffer so that a few may enjoy,
that the one must do the work so that the other may
reap the wages; this unnatural thing, that peoples of one
root and one blood are torn from each other by arbitrary
% --- p. 6 ---
\setcounter{footnote}{0}%
\opage{6}power, while peoples of different descent, with different
language, different customs and different conditions of development,
are to be forced to give up what
constitutes the natural ground of their life, and to live under an
all-levelling uniformity in the same state under
the same laws, only because mighty arbitrariness
wills it.

It is remarkable to see how arbitrariness,
the clerical as well as the secular, sweats and toils
to keep in check the unruliness of the millions and
to secure for civil society order and quiet. It will
not succeed: sin, ah, wicked sin! sin
increases. Faith decreases, corrupted nature grows
worse and worse. Quite naturally; as long as one works
with the principle of arbitrariness, arbitrary reward and
arbitrary punishment, supernatural wrath and supernatural
grace, sin, the unnatural, will surely thrive. Go
to the nature of the matter, give us light on the ground of the laws,
natural relations in the arrangement of society, and the unrest
will cease and sin decrease; then comes the new
time: as things now are, they cannot possibly remain.

The peculiar mixture of the true and the false
which underlies this utterance we shall not here
discuss more closely. Only this shall be stressed, that utterances
of this nature rest on more than the calculus of probability
and statistical tables. They are prophetic
voices we hear; the new prophets come
forward to replace the old.

Every time in which higher vital forces are stirring and
great questions of the spirit are in motion must, in one
form or another, necessarily have its prophets. What
determines the character of an age is the prevailing
% --- p. 7 ---
\setcounter{footnote}{0}%
\opage{7}world-view. Different kinds of world-view may indeed
be contemporaneous and clash with one another;
but, unless everything is in chaotic dissolution, there will always
be fundamental features of a view which has
the future before it. Our ecclesiastical conception of heaven
and earth belongs to antiquity, the scientific one to the present.
The difference between antiquity and the present is in this respect
so thoroughgoing that all signs point toward
a radical transformation. Here is a theme for variations rich
in mood. It belongs to the essence of prophecy
to strike the fundamental chords of existence and touch the strings
that vibrate most deeply in the ear of the generations. As the prophets
stand on eminent heights, they appear as solitary
figures, but if we look more closely, they are by no means
solitary; there are thoughts that keep them company, they have
their life in the midpoint of the forces. What ferments in
the race, its need and longing, its pain and hope,
moves powerfully in them and grows clear to their sight;
the dimly sensed, the not yet explained, that after which the many
grope, what is bound in the mind and has not found the
releasing word, this they fetch out of the night of the deep and lift
up into the light of day.

The relation of the supernatural to the natural, the relation of freedom
to necessity, the will of the ruler of the world
and the connection of events, these immortal
themes of the spirit keep the movement going and constitute
the standing account between old and new prophets.

If we pass over the mythically veiled, in which natural
and supernatural are mingled --- the oracles, the sibyls,
the völvas ---, and keep to the sharply marked
opposition of spirit and nature, then we find
% --- p. 8 ---
\setcounter{footnote}{0}%
\opage{8}what is typical of antiquity in Israel's old prophets.
Jehovah, the God of the spirit, is the creator of the world, his will
almighty. The being and working of the Spirit is therefore not
dependent on finite natural conditions; the almighty
Spirit works when it will, what it will and as it
will, with unconditioned freedom. --- With such an arbitrary
use of the powers of omnipotence the new prophets can in
no way reconcile themselves. A miracle-working will
that does as it pleases with nature and plays
an Olympian game with causes and laws they find to such a
degree contrary to experience and sound reason that,
offended at this, they often go to the opposite extreme
and deny the original freedom of the spirit. It is
then not the free will of the spirit but the immanent necessity of eternal
nature that is the all-determining
first. What nature demands, the spirit must will: the
divine will is, according to this way of seeing, just
as little able to change the laws of nature, the course of nature,
as the human.

This opposition is strengthened still further by
the manner in which the communication of the spirit is supposed to
take place.

The old prophets received the Spirit immediately
as Jehovah's messengers; the Spirit came upon them with
Jehovah's word, proclaiming Jehovah's will. The calling
to be a prophet was a miracle of the Spirit. The relation between
the divine spirit and human consciousness
therefore continued to be miraculous. The prophet
knows how to distinguish between his own thoughts and the thoughts
that are inspired in him by the Spirit, between his own words to
Jehovah and Jehovah's word to him. How circumstantially
and with what clarity the Spirit can foretell to the
% --- p. 9 ---
\setcounter{footnote}{0}%
\opage{9}prophet things which no human being could know of himself,
of this there are manifold examples. When Samuel (1 Sam.
ch.\ 10) had anointed Saul king over Israel, he kissed
him and spoke thus: ``When thou goest from
me today, thou shalt meet two men by Rachel's sepulchre in
the border of Benjamin's lot at Zelzah; they shall say
unto thee: The asses which thou wentest to seek are
found; but now thy father hath put the matter of
the asses out of his head, he is troubled for you and
saith: What shall I do for my son? When thou goest
on forward from thence and comest to the oak of Tabor,
there shall meet thee three men going up to
God to Bethel, one of them carrying three kids,
another three loaves of bread, the third a bottle of wine. They
will salute thee and give thee two loaves, which thou shalt
receive of their hands. After that thou shalt come
to the hill of God, where the pillar of the Philistines stands; when thou
art come thither to the city, thou shalt meet a company of prophets
coming down from the high place with stringed instruments,
tabrets, pipes and harps before them, and they shall prophesy.
Then the Spirit of Jehovah will come mightily upon thee; thou shalt
prophesy with them and be turned into quite another
man than before. And when these signs are come, then
do what thou feelest the urge to do, for God is with thee''.
And it happened exactly so. In the moment Saul
``turned his back to Samuel, God turned his heart, so that
it became another. And all those signs came to pass on
that day. When he came to the hill, a company of prophets
met him, and the Spirit of God came mightily upon him,
and he prophesied among them. All those who
knew him beforetime and saw that he prophesied
with the prophets said
% --- p. 10 ---
\setcounter{footnote}{0}%
\opage{10}among themselves, what is come unto the son of Kish? Is Saul also
among the prophets?''

Such a communication of the spirit, so pettily exact a
prediction of chance encounters can --- according to
the new prophets --- find credence only with those who
feel disposed to every kind of faith, also to faith in
clairvoyance and spiritism. No one who has any
notion of the laws of existence, spiritual as well as
natural, will be able to form any conception of
how a spirit which is not our spirit can step forth
speaking in our inner self and inform us how many
men with kids and bread and bottles of wine we shall
meet on the road. --- If we call the attention of these wise men
to the fact that there are surely many, and that even
among the enlightened, who still sincerely believe in Samuel's
prophetic foresight, they answer us that this happens through
a pious, subtle self-deception. It is with prophetic
images of the future as with natural light-images:
a slight turn of the mirror's surface can change the whole.
Samuel's mind is filled with faith in Jehovah; the spirit
stirs in him at the thought of the decisive step
of having to choose a king for the people. So he speaks to
Saul with an authority which his age and position
and superiority of spirit give him. Saul shall no longer
search for his father's asses; he shall find the prophets
and become a man of the spirit; for only by becoming
quite another man than before ``will he as king be able to
rule Israel''. It works on the gifted young man,
the word strikes down into his soul like lightning; thus we find
him inspired in the circle of the inspired, --- Saul among
the prophets. That is understandable, for it is natural.
When one, half consciously, half unconsciously, substitutes
% --- p. 11 ---
\setcounter{footnote}{0}%
\opage{11}such an explanation, while on the authority of tradition,
and because the supernatural always glitters in
the imagination, one keeps literally to the narrative, or,
as it is so coquettishly and ingratiatingly called, receives the
revealed word in all simplicity, one can of course easily
persuade oneself and others that, even in the age of unbelief,
one is one of the old-fashioned believers. But a cloud of mist
of ambiguities broods at the bottom of such
souls. If self-examination were awakened and the light of
reflection were allowed to break through, the old-fashioned
believers would, surely to their own surprise,
come to make modern discoveries. The letter-believing
conception and the substituted understanding belong
to two utterly conflicting views, that of antiquity and that of
the present.

According to the original meaning, the Spirit of Jehovah
has arbitrarily, i.e.\ without any psychological regard whatsoever
to natural conditions, inspired Samuel with every word in
the prophetic prediction, which is then also confirmed
by the signs' coming to pass punctually. But not
to speak of the fact that so nature-less, arbitrary a mechanics of the spirit
stands in glaring conflict with the inner as well as
outer experiences in which every man of the present breathes
and lives, there is in this at the same time a contradiction which even
the most unconditional faith in miracles is not able to
remove. If Jehovah is omniscient, and the Omniscient
has once, in all arbitrariness, decided to communicate to
a human being what will happen in the nearest future,
then on this presupposition one can find it
in order that the prediction is fulfilled down to the smallest
details. But precisely by holding fast to
the presupposition we are driven by an inflexible logic
% --- p. 12 ---
\setcounter{footnote}{0}%
\opage{12}further on to the question: how is it possible
that the omniscient spirit, which here has hit so exactly
all the side-issues, has missed the main thing? Samuel is in
the spirit assured that Saul is the right man, but
that is, as the following history shows, a great
mistake. Saul is precisely not the right man. Is
it now Jehovah who has erred, or is it Samuel,
or perhaps both of them? However one turns
and twists, one comes into a labyrinth of contradictions.
If the Omniscient knew that Saul really was
the worthiest when he was chosen king, why then did
the Omniscient not also know that he would later make
himself unworthy? If one says: the Omniscient knew very
well that Saul would fail his calling, but chose him
nevertheless, because he would choose him, then we have
pure arbitrariness. If one says: Saul's disobedience was
an expression of his free will, and what will happen in the near
future by virtue of men's free will
not even the Omniscient himself can foreknow; then one comes
into difficulty with the kids, the loaves and
the bottle of wine; unless one will assume that everything was
ordered in advance, and that the men Saul met did not meet
him with what they carried of their free will, but according to higher
orders. If the latter is assumed, what then is Samuel's
whole prediction? A prearranged game, a comedy
of prophecy. To come out of the labyrinth of arbitrarinesses
there is only one way out: to take the natural conditions
along and keep to the substituted
explanation.

With the entry of the natural conditions arbitrariness
departs, and the communication of the spirit appears in another
light. All gifts of the spirit are, originally seen, natural gifts;
% --- p. 13 ---
\setcounter{footnote}{0}%
\opage{13}it is only in the gifted, in those by nature especially endowed,
in a word: in the geniuses, that the great world-thoughts we
call ideas are transfigured into spirit. The ideas do not come
through supernatural, miraculous promptings, still less do
they fall speaking into the inspired soul! The genius
receives his impressions from the surrounding world like other
men; the difference is only that he, by virtue of a
peculiar organic disposition, works up his impressions in
a peculiar manner, re-melts them in memory,
gathers them in intuition, forms them, clarifies them and
translates them into his own peculiar, yet generally intelligible
language.

The inner transformation often takes place in so involuntary
a manner that the gifted man may fall into wonder
at himself and ask: what is it, then, that goes on
in me? It works in here without my working, it
thinks even when I do not think; it grows as
the grass grows, it grows as the fetus in the pregnant woman.
What I bring forth is not of myself;
my activity, even if I exert myself to the
utmost, is quite subordinate, I am only a servant,
a midwife. Thus it may well seem to
him that the new visions with the new thoughts are enigmatic
promptings. But they are not promptings of
a supernatural spirit speaking into the soul; it is
the play of the soul's own powers, the nature of genius, in whose
secret workshop work goes on according to laws. The gifted man
cannot possibly make a separation between what
the enigmatic works in him and what he himself effects.
By attempting such a separation and taking the
natural for something supernatural, he would sink down
beneath a nature-less mechanism of the spirit and cease to be
% --- p. 14 ---
\setcounter{footnote}{0}%
\opage{14}a genius. Genius is precisely the unity of nature and consciousness,
the happy unity in which the activity of nature and
self-activity go together into spirit. The word of genius is
a human word; but this human word, which the thought of nature
has borne and the idea held at its baptism, receives through
its consecration a magic power which wonderfully seizes and
moves our mind. This human word is the word of the
spirit, a prophetic word, with a reflected glory of the divine.

So then: the heralds of antiquity are prophets of supernatural
spirit, the heralds of the present are prophets of genius;
here we stand at the nodal point. Although much of what
the new prophets have to tell us about the laws of nature well
deserves to be taken to heart, there is nonetheless something
which they partly conceal, partly overlook. They maintain the
validity of laws, and we must grant that they are right; where there are no
laws, or where the laws are arbitrarily broken, there is
no connection, no meaning or understanding.
But they speak as if there were no other laws than
the laws of nature, as if the spirit did not also have its
laws. It constantly looks as if the physical
laws of necessity were everywhere to take the place of the laws of the spirit.
Here we must be on our guard; for if the indicated
principle is to be carried through, it has
very far-reaching consequences. For it follows from this not
merely that the divine spirit cannot possibly be
any independent being, distinct from nature, but
it also follows from this that our soul, our self, our will
cannot be any being distinct from our bodily organism
either. Has it then really been proved
that the self-validity of the laws of the spirit cannot possibly subsist with
the laws of nature? By no means! The whole demonstration
% --- p. 15 ---
\setcounter{footnote}{0}%
\opage{15}confines itself to a bare assertion. It must be
admitted that our insight into the laws of the spirit is too imperfect
for us to be able, in particulars, to give a sufficient
account of the interaction between nature
and spirit. But, we ask, can the physics-philosophers,
by which we mean such critical half-thinkers as
live by copying the results of investigating physics,
can they then really, on their side, in all particulars --- by
means of the laws of nature give an account of how the
inorganic becomes something organic, the organism a
living body, the living body something that
looks like a soul, a self-being, a conscious
self-determining will? Can the new prophets perhaps
explain genius? They will surely come to
be brief and let it rest with the categorical
declaration that the extraordinary in genius is subject to the laws
of nature. With the same right we could then assert
that the supernatural in the old prophets'
inspiration is subject to the laws of the spirit. Or has it
perhaps been proved that the supernatural is the
nature-less, the unnatural? No, here is no proof,
here is only an assertion.

If we could not know genius by its productions,
nor appropriate these productions,
rejoice in them and use them, before we had fathomed
to the bottom what genius is, and to which mechanical
or chemical combination of substances it is to be referred,
then we should for the time being have to renounce
art and poetry; we should have to regard the heroes of civilization
as suspect, enigmatic quantities, indeed, what is more,
we would not at all be able to understand the new
prophets. But as surely as we can prize and admire
% --- p. 16 ---
\setcounter{footnote}{0}%
\opage{16}the productions of genius without being able to comprehend the
inspiration of genius, why then should we not acknowledge
the wonderful productions of the spirit of revelation,
although we are unable to explain the supernatural
inspiration. It is not quite as the
most zealous preachers of modern consciousness loudly assure us
in streets and market-places --- especially in the physics-philosophers'
straw-market ---, that the independence of the spirit is
sheer arbitrariness, that this arbitrariness is to blame for
dualism, and dualism to blame for our spiritual misfortune.
It is hardly certain that we shall be rid of the riddles of life
because we are rid of arbitrariness, for the unconditioned
freedom of the spirit is after all something other than lawless arbitrariness;
we dare not, without more ado, rely on death's settling
the account of freedom if we merely deny freedom
and fall quietly asleep in necessity. Necessity is
no more intelligible than freedom. Anyone who feels
a calling to undertake investigations will convince himself
that the problem of necessity is insoluble precisely because
the problem of freedom is insoluble, and conversely; in other
words: that genius in the last instance is just as unfathomable
as the spirit.

Necessity is the unfathomable of nature; the
unfathomable of the spirit is freedom. The original unity
of freedom and necessity is hidden in the eternal
wisdom. Whether we now feel ourselves addressed rather by
the prophets of the spirit or by those of genius, we must trust in
wisdom --- in folly we could not, after all, trust;
we must look at the course of the governance in history. The development
of the human race is a naturally conditioned development of the spirit,
an often hampered, by manifold obstacles crisscrossed,
but on the whole progressing development of freedom.
% --- p. 17 ---
\setcounter{footnote}{0}%
\opage{17}Arbitrariness and the limits of nature belong to the elementary
presuppositions of development. At the limits encountered
there arises a feeling of necessity, in the use of arbitrariness
a feeling of freedom. Life can as
little do without necessity as without freedom; both
rise before the enlarged vision as ideal powers in
ceaseless strife. Shall freedom serve necessity,
or necessity freedom?

Conditions change. Behind the glittering veil of myth
necessity begins a grandiose play which sets
the youthful powers of the race in motion. Mighty
nature takes on human form; it is like
a jest and yet fearful earnest. On the tops of the mountains,
on the bottom of the sea, in the thicket of the forests, in rivers and
trees, everywhere there are human-like beings who
live, feel and will. The struggle of the elements is a struggle for freedom,
and freedom triumphs: the monsters sink, Olympus
is peopled by blessed gods, and from the abyss the sighs of the giants
are heard. It is a time of heroes, a time of geniuses
with fertile thoughts and great deeds, a
spring of civilization, when the faculties shoot up, nature
is consecrated to spirit and the ideas blossom. All is
beauty and harmony, for freedom is necessary and
necessity free. But the glory lasts only a short while;
necessity at play takes off its mask, and what
is revealed? --- The petrifying countenance of fate! Existence
darkens, the gods pale, the courage of freedom sinks; the understanding
spells out the connection of things, and the further
it gets on in the text, the more clearly it shows
itself that freedom was only a dream: real life
is bondage.

% --- p. 18 ---
\setcounter{footnote}{0}%
\opage{18}Conditions change. In order that the necessity of fate
shall not strangle freedom, the eternal wisdom has so ordained
it that the spirit can come into its own. It begins
over again with freedom, but not, as before, in ambiguity.
The freedom of nature was an ambiguity and therefore
a disappointment. The ambiguity is removed as the spirit separates
itself from nature and proclaims itself the Lord of all creation.
The proclamation is in the miracle. The natural suffers,
the miracle breaks through; before the miracle's blow the threads of necessity
must snap and fate sink. It is
the work of liberation: the almighty will of the spirit breathes
upon the powerless wills, sets the goal of freedom before
their eyes and awakens them to renewed struggle. It is the time of
revelation, a time of martyrs: the temporal is sacrificed
for the eternal, the supernatural gains power over
the natural.

Conditions change. The supernatural breakthrough
worked so strongly on the human imagination that
it was disgusted with the natural and was enthusiastic only for
the supernatural. Nature faded; viewed
in the supernatural light, it looked like an old
harlot, of whose youthful beauty hardly a trace was
left. Instead of understanding the declared opposition
of spirit and nature as an expression of truth for a deeper
unity, one sought liberation in an arbitrary breach of
the laws of nature. The real order of things went
under in sheer miracles, legends and holy tales.
Freedom was misjudged because necessity was
misjudged. But conditions change. The eternal wisdom
has so ordained it that nature can anew come
into its own. The time of discoveries, the time of the natural sciences,
% --- p. 19 ---
\setcounter{footnote}{0}%
\opage{19}the time of free thoughts and of free inquiries is
again the time of geniuses, but --- on higher terms.

The opposite conditions go, each according to its direction,
to such strong, one-sided extremes that for that
reason it is difficult for human beings who live in the one
condition rightly to understand what was originally lived
in the other. But we must trust in wisdom.
Only the governing wisdom can give us a key of understanding,
a solution of the riddle, a standard for the judgment
of \emph{old} and \emph{new prophets}.

% --- p. 20 ---
\setcounter{footnote}{0}%
\lecture{Second Lecture.}\opage{20}It is a foreign, long-vanished world, the old kingdom
of Israel, with which we are to make acquaintance. But,
it is objected, this kingdom surely cannot be said to be
foreign to us. The history of the Israelite people is taught, after all,
in practically all our schools and is known even to the
common man. If there is any people of antiquity about whose
life and fate one may presuppose widely spread
knowledge, then it is surely the Jewish one. --- And yet
that kingdom is foreign to us, more foreign than the culture of the Greeks
and the Romans, more foreign than the conditions of primeval times
in the Stone Age, the Bronze Age and the Iron Age,
whose poor remains are dug out of the earth. And why
so foreign? Because that which is to characterize its
innermost essence is incomprehensible to us.

A great change must have taken place, indeed a
transformation, with heaven and earth, since that old
realm of wonders vanished. A transformation? Actual
nature is, after all, in broad outline the same. Old
Egypt with the river Nile, which sinks and rises, exists
still. Mount Sinai still stands, a giant of granite,
unshakable on its broad foot, although the storms of millennia
have whirled about its grey head. And the land of Canaan,
the river Jordan and the Dead Sea, Mount Moriah and
% --- p. 21 ---
\setcounter{footnote}{0}%
\opage{21}the rock of Zion, the heights of Lebanon, the valleys of Esdraelon, indeed the roses
on the meadows of Sharon --- all is here still as in the days of old.

Yes, the natural is here still, but where is the
supernatural? A thunder may roll from the top of Sinai and
the lightnings hiss like serpents; but where is the sound
of the trumpets? Where is the commanding word of the Spirit? The voice of Jehovah
has fallen silent. The storms awake, the tempest still rides
along the old ways, but where is the chariot of the cherubim, and
where is He who once ``rode upon the wings of the storm?''
Where is the arm of the Almighty? The terrible deeds,
the saving deeds, the great signs, when
he stretched out his hand from the cloud? where is it all,
and what has become of it? --- What has
become of it, that I shall tell you, answers modern
criticism, if you really want to know. It has become
what it was and still is and always will be, religious
fantasy, productions of a spirit whose innermost
essence is illusion.

Now if criticism is right, then we must surely be able to
test it. If it is an illusion, then it must, like
every other illusion, vanish when in the light of day
we call the understanding to our aid and apply the measure
of reality. But to whom shall the measure
be applied? Not to the people of Israel in general; for
this people, with all its stiff-neckedness, was nevertheless fickle;
it changed gods, it changed faith, it was moved hither
and thither like the reed and the rush. We must apply the measure
to the outstanding characters, the leaders of the people,
the chosen of the Spirit; we must take their lives and deeds
sharply into view in order, if possible, to discover
the illusion.

% --- p. 22 ---
\setcounter{footnote}{0}%
\opage{22}If a man in our day stepped forth before the
common people and spoke thus: I know who it is that
has stolen your freedom from you and laid the yoke of bondage
on your shoulders. But fear not, take courage!
trust in the omnipotence of the Spirit and trust in me; I shall
break the heavy yoke, I shall, despite obstacles and
dangers, lead you safely to the goal. I need no
worldly means; I use no carnal weapons; for
the will that works in me is not my will, but
his who bears heaven and earth\dots{}, if a man,
as I said, came forward thus, what would we then judge
of him? He is caught in a lamentable illusion.
Possibly he himself believes in his mission and means
it honestly. But he is an enthusiast, a madman,
to put it mildly a fantast. Let him make the attempt and
try his powers; we prophesy a sorry end for him;
our calculus of probability and our statistical tables
prove that he will suffer shipwreck with his attempt and
run aground on the rocks of reality.

But is it not strange nonetheless; in that old
world of wonders the great leaders stepped forth precisely
in such a manner and spoke to the people in such a
language. What are we now to judge? Were those
men, then, enthusiasts and fantasts? This much is certain,
that they did not suffer shipwreck on the probable, but
carried their work through to the end. Judge the supernatural
as one will, here there is nonetheless something inexplicable.
If faith in the supernatural can prove its truth
by giving strength to lift the natural
and to make the exalted actual, then it is not seen how
the supernatural can be the untrue, for faith
in the untrue, an untrue faith, cannot possibly be able to
% --- p. 23 ---
\setcounter{footnote}{0}%
\opage{23}transform an illusion into truth. If faith in
the supernatural is and remains an imagination which at
the very most has only psychological truth, then there must
in that world of wonders have been a time when
psychological truth was stronger than actual
truth; but then it is impossible to understand why the
stronger truth should not, at least as far as that time is concerned,
also count as the higher truth.

They are characters of great psychological truth,
exalted figures, that we have before us. Let us bid them
step forth, that we may rightly contemplate them; we shall name
three: \emph{Moses}, \emph{Samuel}, \emph{Elijah}.

Since Moses is a world-historical character, the
lawgiver of the Jewish people and founder of its constitution,
it is natural that he has been compared
with other great lawgivers of antiquity, with Lycurgus
in Sparta, Solon in Athens and others. But in such comparisons
one usually overlooks the characteristic
difference and makes Moses a sage in the same sense
as Solon. In his essay ``Die Sendung Moses''
Schiller has treated the matter from a rationalist standpoint.


Initiated into the wisdom of the Egyptian priests, exiled
from Pharaoh's court, Moses in the wilderness of Midian at length
worked out how he might best deliver his
people, impress upon it the faith of the patriarchs in the one
true God, and secure its future as God's people by
an expedient legislation. Without having received
any supernatural revelation, he nevertheless, in order
to gain credence with the ignorant people, had to
stoop to a pious fraud, in that he appealed
to a miraculous calling, pretended to stand
% --- p. 24 ---
\setcounter{footnote}{0}%
\opage{24}in a peculiar personal intercourse with Jehovah himself, and
thus stepped forth speaking and acting in the name
of Jehovah.

The unhistorical in such a conception cannot be
proved by a straightforward reference to the testimony of the Old
Testament. These testimonies could, after all, merely
prove either that the pious fraud was carried through so
masterfully that even the foremost among the people
took the pretended miracles, the outward ones as well
as the miracle of the spirit, for full historical truth, or that
the priests, shrewdly sensing the fraud, helped
to inspire the popular fantasy and, in the history recorded thereafter,
perhaps much later, likewise by a pious
fraud, gave the hovering notions of what had happened
a fixed dogmatic-miraculous form and thereby secured
the mode of presentation that is typical of the theocracy.
Here is something that must be openly and unreservedly conceded
to criticism. It cannot be proved that the books of Moses
were written by Moses himself or by an eyewitness.
It cannot be concealed that the intermediate links of conditions
that should make an understanding possible for us are omitted at many
and most important points. It must be conceded
that the believer in miracles, by taking everything according to the letter,
gets not a hair's breadth further, as regards insight into the inner historical connection of the matter,
than the
non-believer who, rejecting the supernatural,
keeps to legends and myths. Both make themselves guilty,
if in opposite ways, of the same error: of
claiming a knowledge of what cannot be known. Our
only guiding star, the only thing we have to hold to,
is the \emph{psychological truth}.

% --- p. 25 ---
\setcounter{footnote}{0}%
\opage{25}That psychological truth cannot coexist
with any known fraud whatever, be it pious
or impious, is surely indisputable. But should it
not be able to coexist with an unconscious self-deception? Even
if we assume this, we should still have gained something that
is certain, namely the recognition that Moses, as regards consciousness of revelation,
stands on the same, essentially
on the very same level as the people of Israel with all its
priests and prophets. The psychological question
thus becomes this: from what spring of life did the men of the Spirit
in the Old Testament themselves believe they drew the
courage, the wisdom and strength which let them overcome
almost superhuman obstacles and enabled them
to carry out their mission?

It is a peculiar sight to see Moses, an unarmed
man, indeed a fugitive, standing over against his majesty
Pharaoh, commanding him in the name of Jehovah and with
the authority of Jehovah: Let my people go! Pharaoh feels
himself the holder of power, and a holder of power does not so
easily let go of his prey. What does he rely on, then, the unarmed
man? In the solitude of the wilderness by Mount Horeb, at the
time when he pondered the possibilities, he was irresolute
and afraid; he felt his powerlessness, and the sigh sounded:
Lord, send another. Now that, in the midst of the stream of events,
he sees reality as it is, looks the
mighty Pharaoh himself in the eye, now he is not at all
afraid. He does not feel his powerlessness; on the contrary, there breathes
a growing self-confidence, an ever-rising feeling of
power in the repeated words: Let my people go! It is
a great change that has taken place in Moses. If
this change should have come about through a self-deception, then
there must be self-deceptions of very different orders;
% --- p. 26 ---
\setcounter{footnote}{0}%
\opage{26}for the facts show that Pharaoh too finds himself
in a self-deception.

Against imaginings and illusion the actual
world has strong means; it is a trial of strength
between Moses and Pharaoh. If it were permissible to place Moses
in line with the geniuses, then in the concurrence
of circumstances under which the natural plagues of the land
come to his aid, and in the shrewdness with which
he knows how to use them and interpret them to his advantage,
one might perhaps find a confirmation of the well-known saying:
genius has fortune on its side. But to conceive Moses as a
genius is to miss the psychological truth. The consciousness of
genius is mythical at its ground; genius in a higher
power is a hero, begotten by a divine father
on a human mother, a demigod in whom the will
is force of nature, and force of nature will: in the genius illusion has
flesh and blood. Moses is, essentially understood,
no genius; the mythically veiled is as remote from his being
as possible. The aesthetics of sensuous appearances
is wholly unknown to him. Nothing of that twilight magic
which lends nature a soul and makes blocks of rock
into manifestations of mighty spirits is to be traced
in him. For his consciousness the earth is earth, the clouds
clouds and the sea sea. He is, in his view of
the phenomena of reality, as sober as a natural scientist
in the nineteenth century, in his judgment
of civil societies so practically sensible that
he would surely gladly have used the statistical tables,
if only he had had them.

Moses, it is said, is a practical genius, a genius of will,
his wonder-working power is his will. But if
the psychological truth is to be sought at this point,
% --- p. 27 ---
\setcounter{footnote}{0}%
\opage{27}where do we then end up? If it is possible that one
who must actually feel that he has all strength and help
to seek in his own will can be caught in the imagination
that it is a will independent of him, a will
outside him, from which the help comes, then we must
admit that such a self-deception comes very close
to the miraculous. How it is psychologically explicable
that a practical genius with a prosaically sensible
view of nature can --- when it does not happen, that is, by
a known fraud, in order to dazzle the people --- feel
inwardly assured that it is not the circumstances
or, in the play of chance, the forces of nature, but Jehovah himself,
who in the hour of danger comes to the aid of his obedient servant,
we are unable to see.

With the enemies behind him and the waters of the sea
before him, Moses and his people are in danger. What does
``the practical genius'' do now? He touches the waters with his
staff. ``He stretched out his hand over the sea'', says the
old Scripture, ``and Jehovah caused the sea to go back by a
strong east wind all that night, and he made the sea
dry land, and the waters were divided''. If we take Jehovah
away and keep the east wind, then chance itself is
miraculous; if we let the practical genius poetize his will
into an almighty Jehovah, then the self-deception is miraculous.
Only on the presupposition that Moses in his inner being
originally felt his will moved and strengthened by
Jehovah's will, and for that reason recognizes the same almighty
will in the strong wind that divides
the waters, is there psychological truth. This psychological
truth is one with the wonder of revelation.

To explain the miracle of the spirit, the psychological wonder,
is impossible. The religious psychology of our age gives
% --- p. 28 ---
\setcounter{footnote}{0}%
\opage{28}only pale, colourless after-images of the original
relation to God. The truth of the inner experiences is a
secret of faith which each must keep to himself.
The critique of faith that is to be applicable to the history
of Moses must rest on the following two principles:
1) all accounts of outward miracles must
take their light from the miracle of the spirit, 2) the psychological
truth of the miracle of the spirit must be recognized by a judgment
of Moses' character and the work that was
his mission.

The enlightened of the present are really at a loss with
the historical element in the many miracles. There are those who
reason thus: all miracles are equally miraculous;
one must therefore either accept them all or reject
them all. Accept them we cannot, without
committing a spiritual suicide; we must therefore reject
them, i.e.\ assign them a place in the great museum of
the legends and myths of antiquity. This way of seeing is only
consistent if one at the same time rejects the miracle of the spirit
and makes Moses' faith an imagination. But
to hold on to the miracle of the spirit and place Moses on the level of revelation
among the exemplars of faith, although one
leaves his other miracles standing in all their ambiguity,
will in no way do. If the notion Moses
himself has of a miracle is ambiguous, then his notion
of Jehovah is surely ambiguous too, and his faith ambiguous.
Moses is then a man of an ambiguous character.
If on the other hand all is right with Moses' own faith; if there is
no ambiguity, but full clarity, in his inner being
when he speaks of his mission and of all that
Jehovah himself reveals to him, and thus no ambiguity
in his notion of the miracle of the spirit, then neither
% --- p. 29 ---
\setcounter{footnote}{0}%
\opage{29}can there be any ambiguity in his notion of
the miracles of events in the outward world. If Moses is entitled
to come forward as Jehovah's envoy, then
he is also entitled to regard every extraordinary
event that serves his purpose as a helping hand
from Jehovah, a sign from the Almighty.

Here there is no question at all whether Moses
was able to draw any objective boundary between
what Jehovah works through natural causes and what
happens in a supernatural way. That the waters of the
Red Sea were divided according to Jehovah's will, of this Moses
is convinced, but the physical-critical-metaphysical relation
between the east wind and the divine will
he leaves to the modern consciousness to discover. --- So
then, it is objected, the difference between miracles
and natural occurrences can from this standpoint be both
posited arbitrarily and effaced arbitrarily; religious
fantasy thus has just as free a play in the doctrine of revelation
as in mythology. Where is the critique of faith now,
and what does one get out of the miracle of revelation?
Not to speak of how naturally it comes to the peoples
of antiquity to introduce the gods speaking and acting,
it is surely well known that enthusiasts and fanatics have
pretended to inner revelations without its being possible to
conclude from this that they were deceivers. Supernatural visions can
surely, as physiology shows, be explained as hallucinations,
natural outbursts of a religious delirium.

We would point out that it is not a doubling of imagination but
a doubling of will, an inner opposition of divine
and human will, about whose personal
character the inquiry turns. By the fruits
shall the tree be known, by the actions shall the will be known;
% --- p. 30 ---
\setcounter{footnote}{0}%
\opage{30}by the clear grasp of life's oppositions shall
the understanding be known. We ask: is there in the character
of Moses a single trait which points to a mythical intermingling
of the divine and the human? Does
his inspiration betray anything that could arouse suspicion of
ecstatic confusion? Do the spiritual laws from Sinai perhaps bear
traces of a sick brain, a disquieting stamp of
hallucinations?

Can physiology give an account of how a
man whose innermost being is imagination and error
is able to carry out a giant's work such as this: to save
a people from bondage? Is there any psychology that
explains how it comes about that a human being first
converts his own will into the form of an almighty will of God,
and then, without suspecting the monstrous self-deception,
bows under this will of God as an obedient servant,
labours with an ungrateful people through a long, long life,
without weakening, without growing dull, is there
really any psychology that explains this to us?

``And Moses went up from the plains of Moab unto the mountain
of Nebo, to the top of Pisgah, that is over against Jericho,
and Jehovah shewed him the land of Gilead unto Dan; and
all Naphtali, and the land of Ephraim, and Manasseh, and all
the land of Judah unto the sea behind it; and the south, and
the plain, the valley of Jericho, the city of palm trees, unto Zoar. And
Jehovah said: This is the land which I sware unto
Abraham, unto Isaac, and unto Jacob,\dots{} but thou shalt not go
over thither\dots{} And Moses was an hundred and twenty years old
when he died; his eyes were not dim, nor did
his strength depart from him. But there arose not a prophet
since in Israel like unto Moses, for with him Jehovah had
stood face to face.'' (Deuteronomy ch.\ 34.).

% --- p. 31 ---
\setcounter{footnote}{0}%
\opage{31}The next prophet we named was Samuel, the
last among the judges of Israel. His time is another,
his more limited sphere of activity another, but
the spirit is essentially the same. Early, already as a
child, Samuel was called to be a prophet. There rests over
his calling an indescribable glow of idyllic peace
and yet of spiritual, tense, presentient expectation. Dedicated
by his pious mother to the service of Jehovah, he is sent
as a little boy to Shiloh, to the high priest Eli. Eli
is weak and his sons dissolute; it bodes ill for
the house of Eli. Samuel sleeps near the sanctuary where
the lamp of Jehovah burns. He sleeps securely in the still
night, but wakes, for a voice calls. He
believes that it is the old man, and goes in to him and
asks: Lord, didst thou call me? When it happens
again, Eli understands that it is Jehovah who is calling
the boy. ``Therefore Eli said unto Samuel: Go,
lie down! and if Jehovah call thee, thou shalt say:
``Speak, Jehovah! for thy servant heareth.'' Then comes the revelation.''
Jehovah said to Samuel: Now I will do
a thing in Israel at which both the ears of him
that heareth it shall ring\dots{} I will judge the house of Eli for ever
for the iniquity that he knew how
his sons brought a curse upon themselves, and yet he rebuked
them not.'' What is Samuel to do now: a child
and yet a prophet, he is afraid to speak to Eli of
the vision. But Eli called Samuel and said: ``What
was it that he spake unto thee? Now thou must not hide
it from me. God punish thee in his wrath if thou
hidest anything from me of all that he spake unto thee.
Then Samuel told him every whit, and hid
nothing from him.'' --- Wonderful is the seeming
% --- p. 32 ---
\setcounter{footnote}{0}%
\opage{32}composure with which the old high priest hears his
sentence. He does not repent, for his soul is faint; he
does not act, for his will is lame, he does not pray,
for his hope is extinguished. With an almost apathetic
resignation he answers: ``He is Jehovah and may do
what seemeth him good.''

And the omens were fulfilled. Samuel grows up
to manhood and is honoured by the people; as judge he upholds
the law of Jehovah and in his old age receives the peculiar
calling to choose and anoint the first two kings in
Israel.

In the history of Samuel the supernatural is so peculiarly
interspersed in the natural that understanding
seems near at hand and yet suddenly eludes us when we
think we have it in hand. It is the psychological element
that attracts and repels. Samuel finds himself at a
critical turning point, for the Mosaic theocracy is
about to change its form. The time of the judges was a
time of unrest and lawlessness, ``every man did
what seemed good to him''. No wonder the people found
the state of things intolerable and desired a king. Although
it was a demand that in no way conflicted with
the law, it is nevertheless understandable that it had to arouse
concern in Samuel. Not to speak of the fact that the dignity of judge
was thereby lost to his family, there was
here a religious moment of special weight and significance.
Had the people's demand sprung from faith in Jehovah,
the eternal King; was it not rather from worldly desire
and self-wise considerations? Where was the man to be found who,
clothed with earthly royal power, would unconditionally be
Jehovah's servant? It is understandable that the old
prophet, partly from hierarchical jealousy, but chiefly from zeal
% --- p. 33 ---
\setcounter{footnote}{0}%
\opage{33}for the right of the Spirit against encroachment on the king's part, may
have found it necessary to keep a sharp eye on
Saul and, given occasion, to put him to the test.
The Philistines took the field with a terrible army.
``Saul was yet in Gilgal, and all the people that were with
him were terrified. He tarried seven days, until the
time that Samuel had said, but Samuel came not to
Gilgal, and the people were scattered\dots{} Then Saul said: Bring
hither to me the burnt offerings and the thank offerings, and he offered
the burnt offering. As soon as he had made an end of the burnt offering,
Samuel came, and Saul went out to meet him, to
salute him. Samuel said, what hast thou done? And Saul
said: When I saw that the people were scattered from me,
and that thou camest not after the appointed number of days, and
the Philistines gathered themselves together at Michmash, then I thought: Now
will the Philistines come down upon me to Gilgal, before
I have appeased the face of Jehovah; therefore I made bold
to offer the burnt offering. And Samuel said to
Saul: Thou hast done foolishly,\dots{} now thy kingdom shall
not continue. Jehovah will seek him a man after
his own heart and appoint him prince over his people,
because thou hast not kept that which Jehovah commanded thee.''
It is, finally, understandable on grounds of general psychology
that Saul, who is rejected by the prophet, and who
more and more comes to suspect that it is David who
is chosen to be his successor, becomes bitter and
dark, and sees treachery and intrigues in those
nearest him.

No one shall say that in this history the spirit has
encroached too closely on nature. Saul and David are not shadowy
beings, theocratic abstractions; they are
individualities in flesh and blood with great powers and
% --- p. 34 ---
\setcounter{footnote}{0}%
\opage{34}great passions; what contrasts of light and shadow!
Saul's demon and David's harp-playing, the wild outbursts of
Saul's hatred and David's wonderful good fortune: the spear is hurled,
but it does not strike him. Battles and adventures,
friendship, enmity, love and poetry, all unite
to make this first strife over the throne in Israel
as attractive as possible. In Gibeah under a tamarisk
Saul sits on a height with his spear; dark is
his countenance, but his spirit is darker still. His
servants, who stand about him, must hear from his
mouth that they are traitors all together. What pain,
that Jonathan himself is in league with the son of Jesse! Ahimelech
in Nob is called to account, and the sword of Doeg the Edomite
cuts down the priests of Jehovah. But Saul is doomed to destruction;
his strength is broken. At last he is only a
ruin of a soul, a burnt-out volcano, a despairing man who
under cover of night seeks counsel of the dead. He has recourse
to the woman at Endor. She excites his fantasy and
conjures shades up out of the grave. The dead man speaks: ``Wherefore
wilt thou ask of me? Jehovah is departed from thee and is become
thine enemy.'' --- The prophet's voice! Saul is condemned and
sinks down powerless before Samuel's shade.

But just as the consciousness of the present is, along this way,
on the point of appropriating the content of the history, there suddenly
appears the repellent element: the will of Jehovah with the
theocratic mechanism breaks off understanding. ``What
arbitrariness, to make it a crime of Saul's
that, in the extreme moment of danger, he will as king provide
for the defence of his country, instead of wasting time
waiting for a man who delays and does not keep his
word''. Yet from the standpoint of revelation, that is, of the miracle of the spirit,
Samuel is right. What is demanded of Saul
% --- p. 35 ---
\setcounter{footnote}{0}%
\opage{35}is unconditional obedience to Jehovah's law, is in other
words: spirit-will, not self-willed drive-will. By
mistrusting the prophet Saul has himself undermined the ground
on which he was to stand, for only by virtue of the prophet's
choice is he king. We shall admit that the miraculous
relation to God on which the prophetic authority rests
cannot become transparent to us; but we shall
also recognize that this relation to God is neither speculated up
nor reasoned up, but originally posited by
the Spirit as spirit; we shall recognize that the people of Israel
had its spiritual life in this relation with just
as fully valid truth as any people of the present can have
its worldly life in an idea of culture; we shall recognize
that what we call the harsh, the arbitrary, the
mechanical, is precisely the unambiguous expression of the
separation of God's will and man's will which conditions
an actual manifestation of spirit-will.

We now go further forward, to meet Elijah
the Tishbite. Of this prophet too it holds that, although
the conditions of the time have changed, the spirit is nevertheless
the same. After having torn itself loose from the kingdom of Judah
and having had the calf-worship introduced in Dan and Bethel,
Israel is in full course of changing the calf-worship
into the grossest idolatry. It then comes to
a violent struggle between Jehovah and the idols. Omri's
son Ahab, who has taken Jezebel, a daughter of King
Ethbaal of Sidon, to wife and built Baal a house of prayer
in Samaria, is a zealous idolater. There is talk of four
hundred and fifty prophets of Baal and of four hundred
prophets of Astarte, who ate at Jezebel's table.
Against this abomination, ruinous to the people, the prophet
Elijah comes forward. His will is a holy
% --- p. 36 ---
\setcounter{footnote}{0}%
\opage{36}passion, the spirit in him is fire and flame; in his zeal
he ventures all, he sets heaven and earth in motion.
He speaks to Ahab: ``As Jehovah, the God of Israel,
liveth, whom I serve, there shall not be dew nor rain these years,
but according to my word.'' Then he flees,
not from cowardice, but at Jehovah's command; he drinks
of the brook Cherith, and the ravens bring him food. When
the brook is dried up, he goes to the widow in Zarephath;
the meal in the barrel does not waste, and the oil in the cruse
does not fail: he raises her son from the dead.
When the plague on the land had lasted into the third year, and the famine
had risen high in Samaria, Elijah received the command to
show himself to Ahab. Ahab meets him with these words:
``Is it thou that troublest Israel?'' The prophet answers:
``it is not I that trouble Israel, but thou it is,
and thy father's house, because ye forsake the commandment of Jehovah and thou
followest the Baalim.'' It now comes to a decisive
trial. All Israel with the prophets of Baal and Astarte
are gathered, at Ahab's command, about the prophet of Jehovah on
Mount Carmel. ``Then Elijah came unto all the people
and said: How long will ye halt on both sides?
If Jehovah be God, follow him; if Baal, then follow
him! But the people answered him not a word.'' Only
when the trial of sacrifice was ended, and it had been shown that the
deaf idol gave no sound, however much his priests
cried and cut themselves till the blood came with knives, but that
Jehovah answered Elijah by the fire that fell down and
consumed his burnt offering: then the people threw themselves on
their faces and cried: ``Jehovah is God, Jehovah is God!''
The prophets of Baal are slain, and the plague on the land ceases.
But Jezebel rages. Elijah must flee to the wilderness; despondent
and weary, he wishes for death; ``Jehovah, now take
% --- p. 37 ---
\setcounter{footnote}{0}%
\opage{37}my life, for I am not better than my fathers.''
He is strengthened and refreshed and comes to Horeb,
the mount of God; zeal will consume him, he complains and
prays: ``God of hosts! The children of Israel have forsaken
thy covenant, thrown down thine altars and slain thy prophets
with the sword; I alone am left, but me too
they seek to take my life.'' But the cool
winds cool him, a breath of holy devotion brings
calm to his mind and subdues his passion. Jehovah
himself will come. ``A great and strong storm, which rent
the mountains and brake the rocks, went before, but Jehovah
was not in the storm. After the storm came an earthquake,
but Jehovah was not in the earthquake. After
the earthquake came lightning, but Jehovah was not in
the lightning. After the lightning came a still, soft sound;
then the prophet wrapped his face in his mantle'', and perceived
in all meekness the presence of Jehovah.

Only a few more steps, and he is at the goal.
After having called Elisha to be his disciple, and foretold
to Ahab and Jezebel a bloody, ignominious death, he ends
his stormy life in a tempest. ``He went away
from Gilgal, and said to Elisha: Tarry thou here, for Jehovah
sendeth me to Bethel; but Elisha answered: As surely
as Jehovah liveth and thy soul, I will not leave thee.
So they went in haste from Gilgal to Bethel, from Bethel to
Jericho, from Jericho to Jordan, where Elijah wrapped his
mantle together and smote the waters, so that they were divided.
As they went over, Elijah said to Elisha: ``Ask something
that I shall do for thee'', and Elisha prayed: ``Might
there but be a double portion of thy spirit upon me!'' The answer
sounded: ``Thou hast asked boldly. If thou see me taken
% --- p. 38 ---
\setcounter{footnote}{0}%
\opage{38}away from thee, thou shalt have it; if not, then thou shalt have
nothing.''

How foreign it is to us and remote, that old
kingdom of wonders of Israel! What use is it, then, to pretend
that one literally believes the history of Elijah, and then,
when it comes to an explanation, to have nothing else
to rely on than some elements of general religious
psychology with phrases and assurances and all manner of
probabilities. Only the Spirit itself can recognize what
the Spirit wondrously brings forth. Had Elisha merely seen
with sensuous eyes, he would no doubt have seen lightning and
storm clouds, but not Elijah as he was taken
from him; had he only seen the actual in a natural
light, he would not have seen the ascension. But
that in the midst of the sensuous he really saw the supersensuous,
and inwardly in the spirit of Jehovah experienced the miracle,
we hear from his outcry: ``My father, my father, the chariot of Israel
and the horsemen thereof!''

The actuality of revelation is actuality of spirit, actuality
for the Spirit.

% --- p. 39 ---
\setcounter{footnote}{0}%
\lecture{Third Lecture.}\opage{39}The authors at the Old Testament stage stand
so immediately under the impression of the supernatural that,
in the re-forming activity in which the accounts
were written down, they typically let the miraculous break
in with abrupt transitions and surprising particulars.
The events are not photographed; they have
passed through a re-forming in religious imagination.
As far as the principle goes, one may very well remain
convinced of the psychological validity of the miracle of the spirit
without thereby inferring the literal infallibility
of the historical tradition. The interpreters at the present
stage can have only a mediate and derived impression
of the supernatural; they are referred to a conception
that places the natural conditions of the soul in the foreground;
they must content themselves with a reflected gleam of the
miraculous, an afterglow, which sets its divine
stamp on the thinking, willing and acting of the great characters.
In deep acknowledgment of the wonderful strength and the
soul-powerful truth of the prophetic characters,
the psychology of faith has at all times a religious element
in which it can meet with the psychology of revelation.


% --- p. 40 ---
\setcounter{footnote}{0}%
\opage{40}When the prophetic spirit is to be determined in particular
as a spirit of prophecy, let one beware of
conceiving prophecy without spirit. Some examples.

\emph{The blind Ahijah} and \emph{Jeroboam's wife}
(1 Kings ch.~14). The same prophet who had foretold to
Jeroboam that he should become king over the ten
tribes comes in his old age to foretell to him
that his downfall is near, that his house shall fall.
The closer circumstances are well known: Jeroboam's
son Abijah fell sick; the king, who knew in himself
how deeply he had offended the prophet, has his
wife disguise herself and go to Shiloh, to hear a
prophecy of the son's life. ``And Jeroboam's
wife did so, and arose, and went to Shiloh, and came to
the house of Ahijah. But he could not see, for his eyes'
light was quenched by age.'' If we now keep the
supernatural in the twilight of the background, then it is not
inexplicable that the blind seer may have recognised
the disguised queen, that he may have foretold
death and ruin over the whole idolatrous house, and
that the image of fall and desolation which the old
blind man saw in the spirit may in its particulars have
been all the more terrible the stronger his inspiration
was at the given moment. This is indeed also the narrative's
fundamental character in terms of general psychology. But the psychology of revelation
demands a striking, miraculous working
of the supernatural, and the account runs: ``Jehovah
said unto Ahijah, Behold, the wife of Jeroboam cometh
to ask a thing of thee for her son; for he is
sick: thus and thus shalt thou say unto her: for when she
cometh in, she will feign herself to be another woman. And as
soon as Ahijah heard the sound of her feet approaching the door,
% --- p. 41 ---
\setcounter{footnote}{0}%
\opage{41}he said, Come in, thou wife of Jeroboam; why feignest
thou thyself to be another? for I am sent to thee
with heavy tidings. Go, tell Jeroboam, Thus
saith Jehovah, the God of Israel: Forasmuch as I exalted thee over
the people to which thou didst belong, and made thee prince over
my people Israel, and rent the kingdom away from the house of David, and
gave it thee; and because thou hast not been as my servant
David, but hast done more wickedly than all\dots{}, behold,
therefore I will bring evil upon thy house\dots{} Him
that dieth of Jeroboam in the city shall the dogs eat;
and him that dieth in the field shall the fowls of the air devour;
for Jehovah hath spoken it. Arise thou therefore, get thee home,
when thy foot entereth into the street of the city, the child shall die.''

It will be seen that it can be neither the literal
holding fast to all particulars, nor the
naturalistic explaining away of the supernatural, still
less the limping mediation between the probable and the
improbable, in which an interpretation is to be sought here.
Anyone unbiased must be able to see that it is precisely the refraction
of supernatural and natural light in the prism of the psychology of revelation
that gives the picture its character
and colour.

\emph{Nebuchadnezzar's dream} (Daniel ch.~2). In the psychology of revelation
the gaze is directed immediately toward
the miraculous, whether it comes chiefly from within
or from without. By what, then, shall a true prophet be
known from the false? By the fact that the
Omniscient reveals to him things of which no human being
can of himself have any knowledge. It is this proof
that is brought before Nebuchadnezzar, the king of Babylon,
the first destroyer of Jehovah's temple. In the second year of his reign
Nebuchadnezzar dreamed a dream; ``and his spirit
% --- p. 42 ---
\setcounter{footnote}{0}%
\opage{42}was troubled, and his sleep brake from him. Then the king
commanded to call the scribes, and the enchanters, and
the sorcerers, and the Chaldeans, for to shew the king
his dream; and they came and stood before the king.
And the king said unto them, I have dreamed a dream, and
my spirit was troubled, so I will know the dream. Then spake the Chaldeans
to the king in Aramaic, O king, live
for ever: tell thy servants the dream, and we will shew
the interpretation. The king answered and said to the Chaldeans,
The command is gone forth from me: if ye will not
make known unto me the dream, with the interpretation thereof, ye
shall be cut in pieces, and your houses shall be made a
dunghill.'' Nebuchadnezzar is at one and the same time a wise
man and a mad man, a practical genius and a visionary.
The Chaldeans beat about the bush; but the king stands
firm on this: as surely as they can \emph{interpret} the dream to him,
they must also be able to \emph{tell} him the dream. If they cannot
do the latter, then the former is lies and deceit.
An unbending and sound logic! On this logic
the Chaldean wise men were executed. Daniel
is called; he is brought before the king and says: ``The secret
which the king demandeth to know cannot the wise men,
the enchanters, the scribes, the astrologers shew
unto the king. But there is a God in heaven that
revealeth secrets, and he hath made known to
Nebuchadnezzar what shall be in the latter days.'' In the
revelation of secrets the miracle of the spirit seems
all at once to become so unconditional that every glimmer of
natural understanding must vanish; the religious psychology of the present
is close to despair. And yet there are
hints here of a general-psychological nature. Facing the
mighty conqueror, Nebuchadnezzar drunk with his own greatness and power,
% --- p. 43 ---
\setcounter{footnote}{0}%
\opage{43}stands Daniel with a composure, a clarity of survey
and a superiority of spirit which recall
Moses when he stood before Pharaoh. --- What
does a Nebuchadnezzar dream of? Of his own greatness and that of his
realm! What sort of dream-images can they be
that, in visions over his head about his
bed, could cause his spirit to be troubled? Images of
greatness and of the fall of greatness! ``Thou, O king, sawest,
and behold a great image. This great image, whose brightness
was excellent, stood before thee; and the form thereof
was terrible.'' Straightway at this beginning
the prophet can read in the king's countenance that he has hit upon
the right thing; he continues: ``This image's head
was of fine gold, his breast and his arms of silver, his belly
and his thighs of brass, his legs thereon of
iron, his feet part of iron and part of clay.'' The
psychological stress lies on this: that through Daniel's words the content of the dream
comes into being anew before Nebuchadnezzar's
inner sight, and surely in clearer, more definite,
more significant features. ``Thou sawest till that a stone
was cut out without hands, which
smote the image upon his feet that were of iron and clay,
and brake them to pieces. Then was the iron,
the clay, the brass, the silver, and the gold, broken to pieces together, and became like
the chaff of the summer threshingfloors; and the wind carried
them away, that no place was found for them: and
the stone that smote the image became a great mountain,
and filled the whole earth.'' Nebuchadnezzar listens and
ponders; he dreams his dream over again as he hears
the prophet's interpretation; and the interpretation is in truth
magnificent.

% --- p. 44 ---
\setcounter{footnote}{0}%
\opage{44}The conception can now swing between two poles, the
miraculous and the general-psychological, but always
within the sphere of the spirit. If there is truth in the dream
and its interpretation, then there is also reality; but
it is not political reality, with statistical tables
as the guiding thread in the history of states; it is, with all its historical
weight, reality of the spirit, reality for
the spirit.

\emph{Messianic prophecies}. As long as the miracle of the spirit
is in full psychological activity, and faith in revelation
lives in its original relation to God, the
prophecies can be taken singly, indeed sentence by sentence, without
the conception thereby becoming spiritless: a name, an image,
an ambiguous expression is often enough to set a stream
of spiritual life in motion. But when the impression of
the original has grown faint, when the present has
become so poor in spirit that it lives only in shadowy
re-formings of the past and must console
itself with the hope of something to come, then sentence-by-sentence
interpretations of old prophecies are a doubtful
matter.

The prophets have prophesied of a Messiah who
was to come. Can one really show that these
prophecies fit Jesus of Nazareth and no
other? The Christians assert it, the Jews deny it,
both parties with equal right, so long as one keeps one-sidedly
to words and sentences. It cannot be denied from this
point of view that the objections of the rabbis
have often shone with sharp criticism and obscured the defence of our
ecclesiastical dogmaticians. We choose a single
passage (Isa.~7:14--16): the prophet speaks to the king of Judah,
who, threatened by overpowering enemies, in his despondency does not
% --- p. 45 ---
\setcounter{footnote}{0}%
\opage{45}dare to ask a sign of salvation from Jehovah: ``Therefore shall
Jehovah give you a sign: Behold, the virgin shall conceive
and bear a son, and thou shalt call his name
Immanuel; milk and honey shall he eat, until he
learns to refuse the evil and choose the good. For
before the child learns to refuse the evil and choose
the good, the land shall be made desolate before whose two
kings'' (Rezin, king of Syria, and Pekah, king of Israel)
``thou art in dread.'' This, now, according to orthodox
dogmatics, is absolutely to be a messianic passage of Scripture,
indeed a principal passage. ``The virgin-born, Immanuel,
the God-man, a sign of salvation! --- a passage of Scripture which
is moreover cited by the evangelist Matthew (ch.~1,
v.~22--23): he who cannot see how the prophet
here has \emph{beheld} in the \emph{spirit} the great saving wonder,
that Jesus in the fulness of time should be born of Mary, the pure
virgin, must not have theological eyes in his head.''
There is needed, however, neither rabbinical learning nor
super-theological acuteness to discover that the
orthodox assertion rests on weak grounds. Of a
miraculous virgin birth there is no question here at all. The prophet
speaks, according to the words of the original text, not of the
pure virgin, but rather of his young wife, whose
sons (cf.~ch.~8, v.~3) he gives symbolic names.
The salvation of which he gives warning is so far from being
placed in a distant future that, on the contrary, by the
words ``before the child learns to refuse the evil
and choose the good'' it is expressly made known how
near the salvation is. If the land is to be saved before
Immanuel has learned to distinguish between good and evil,
how then can Immanuel be the Saviour? When,
on the other hand, the evangelist, using the Alexandrian
% --- p. 46 ---
\setcounter{footnote}{0}%
\opage{46}translation, looks back on the passage in Isaiah, stresses
the word ``the virgin'', and so finds the prophecy
fulfilled a second time, indeed in a higher sense perfectly
fulfilled, then in the name of revelation we must grant that he is
right; but his right is in the spirit.

Regarded in the light of the spirit, the whole Old Testament
is a prophecy of the New, as surely as the growing
youth is a prophecy of the man, the unfolding
flower a prophecy of the fruit. The law is
the foundation of the Gospel, and the Gospel is the fulness of the law.
But the transition does not take place quietly and organically as in the life of plants;
the higher stage of revelation is not related
straightforwardly to the lower as the deed of manhood to the
dreams of youth.

What gives the law of Moses the stamp of divine
revelation, namely the unconditional presupposition that
every one of its commandments is immediately dictated by Jehovah
himself, enables the spirit-will to assert its sovereign right
over the natural side of human life in all civil and moral
relations, and to found a state that allows
political life no particular independence
over against the religious. A consequence of this was that the prophets,
although not invested with worldly power,
could exercise a mighty influence on the course of
government, and their authoritative word be decisive for war
and alliance with nearer and more distant peoples. It
is related (Isa. ch.~7) that when the rumour of the threatening advance of Rezin and Pekah
spread in Judah, ``then the
king's heart trembled, and the heart of his people, as the trees of
the wood tremble before the wind''. At Jehovah's command
Isaiah goes to meet Ahaz and says: ``Take heed, be
quiet, fear not, neither let thine heart be faint\dots{}
% --- p. 47 ---
\setcounter{footnote}{0}%
\opage{47}for these two smoking firebrands, for the snorting wrath of Rezin and of Syria and of
the son of Remaliah. Syria hath taken evil
counsel against thee, with Ephraim and the son of Remaliah, saying:
Let us go up against Judah, and vex it, and take it to
us, and set Tabeal as king over it. But thus saith
the Lord Jehovah: It shall not stand, neither shall it
come to pass.'' Ahaz and the people see the danger where it is not,
and do not see it where it is. They would make
an alliance with Assyria, perhaps even seek help from
Egypt, and do not consider that when these great heathen peoples
gain entry into the land, they will trample everything under
foot and destroy the theocracy. Egyptians and Assyrians
do not come as helpers; they come only on the day of wrath
as Jehovah's instruments to carry out a sentence of punishment.
``In that day shall Jehovah hiss for the fly that
is at the borders of the rivers of Egypt, and for the bee that
is in the land of Assyria; and they shall come and encamp together
in the shut-in valleys, in the clefts of the rocks, and upon
all thorn hedges, and upon all meadows \dots{}'' In every
place where there were a thousand vines at a thousand silverlings'',
oxen and sheep and wild beasts now graze.
``With arrows and with bows shall men go thither, for
all the land shall become thorns and thistles.''

Humanly regarded, the life of the prophets is unceasingly in
danger. The authority of the word and the power of kings, the right of
eternity and the demands of the moment, the requirements of the law and
the passions of the people: here are unequal weapons! The people of
Israel, which according to its fundamental determination is a people of the spirit,
Jehovah's people, is according to its elementary natural determination
at the same time a sensual, selfish and hard-hearted
people, a God-fearing and yet ungodly people, faithful and
yet faithless, obedient and yet defiant, courageous and yet
% --- p. 48 ---
\setcounter{footnote}{0}%
\opage{48}faint-hearted, self-sacrificing and yet egoistic, a contentious and contrary
people, at strife with its God and at strife with
itself.

The clearest testimony to the power of the spirit in
prophetic inspiration is the honesty, the purity of will
and the unshakeable courage with which Jehovah's
messengers carry out his errand. The prophets are men of the people
in the true sense of the word; they live with
the people and for the people, they comfort it in the time of trial,
they strengthen its hope, they awaken its strength and open
bright prospects to it; but always on one condition, the
unalterable condition that the people will honestly keep
to the word and spirit of the law. Never do we hear them say:
this people is a people of God, for its heart is as
God's heart; what such a people of God with a God's
heart has now got a vision of, even if it were
a little calf-worship, must surely be the right thing: the people
knows it, the people understands it. No, when the people
sins and the spirit is wrathful, then the word sounds stern; we
hear Isaiah (ch.~1) cry: ``Woe, a sinful people! ---
The ox knoweth his owner, and the ass his master's crib;
but Israel doth not know, my people doth not
consider.''

The distinguishing mark of the revelation of the law must be sought precisely
in this: that the spirit-will stands unveiled, pure and clear over
against the drive-will; the sensuous self is to bow
unconditionally to the command of the spirit. That was the hard thing in
Israel's election, that the divine spirit came
so immediately near to the natural man. Against
the wrath in the prophet's word the heart of the people was hardened,
for truth wounds; when the spirit shone like the sun of heaven,
the eye was blinded by the strong rays.
% --- p. 49 ---
\setcounter{footnote}{0}%
\opage{49}Truth and hardening, revelation and blinding
correspond to each other psychologically. The blinded people
slew its prophets; the relation to spirit made
the children of Israel criminals against the spirit: what a wondrous
union of blessing and curse! It was the condition
of freedom, and yet a fate so hard that love
itself had to weep over it. ``O Jerusalem, Jerusalem, thou
that killest the prophets, and stonest them which are sent
unto thee, how often would I have gathered thy children together, even
as a hen gathereth her chickens under her
wings, and ye would not!''

That Jehovah's law had to be fulfilled, on that prophets
and priests were unconditionally agreed. But fulfilment
has two sides, an outward and an inward. There is
an observance of the single commandments in the outward which
is to a certain degree conditioned by obedience, but the obedience
is only a semblance, for the disposition does not answer to
the deed. The fulfilment which the law of the spirit demands
is essentially inward, an obedience in mind and heart,
an observance in deed and in truth. By casting a
glance at the unequal conditions of the priests and the prophets in
the Jewish theocracy it will become intelligible to us that
the former kept chiefly to the outward, the latter above
all to the inward. The priests, the guardians of the holy tradition,
the supports and bearers of the established order,
had the calling to attend to the sacrificial service in the temple,
to read out the law, to provide for purifications, sabbaths and feasts,
and to see closely that the religious ceremonies took place
in the prescribed order. Their activity
was what one calls official, and so too was their
position. The tribe of Levi, with fixed cities and fixed revenues,
tithes, prescribed vestments, etc., \emph{furnished}
% --- p. 50 ---
\setcounter{footnote}{0}%
\opage{50}the officials of the theocracy. The prophets were not
officials, had neither tithes nor vestments
nor fixed dwellings; they lived freely as the birds under
heaven, and made do with the uncertain: if men
would not provide for them, then the ravens had to do
it. --- It was with solemn ceremonies that Moses
(Lev. ch.~8) consecrated Aaron. ``And he put upon him the
coat, and girded him with the girdle, and
clothed him with the robe; and he put the ephod
over it, and he girded him with the curious girdle of the ephod.
And he put the breastplate upon him,
and in the breastplate he put the Urim and the Thummim. And
he put a high mitre upon his head, and upon the mitre
in front a golden plate, the holy ornament.'' So was
the high priest Aaron in full vestments. In like
manner the priests, Aaron's sons, were furnished with
coats and girdles and high mitres. Here was something
outward, something for the assembled congregation to look
at. With the calling of the prophets it was
otherwise. ``Elijah found Elisha the son of Shaphat, who was plowing
with twelve yoke of oxen before him (he himself was with the
twelfth); as Elijah passed by he cast his mantle
upon him (1 Kings ch.~19); with that Elisha was called.
The spirit makes no great ceremony; it comes, and
then it is enough: Amos, the herdsman of Tekoa, was called
from the flock of sheep.

It is intelligible that the prophets made a quite
different demand upon the inward than the priests did. The demand,
the necessity of an inward fulfilment of the law,
Isaiah (ch.~1) has uttered with a voice that sounds
through the ages, and with an emphasis of truth
that can make outwardness itself tremble: ``Hear
% --- p. 51 ---
\setcounter{footnote}{0}%
\opage{51}the word of Jehovah, ye princes of Sodom; mark the word of God, ye
people of Gomorrah! To what purpose are your many
sacrifices unto me, saith Jehovah; I am full of the burnt offerings of rams
and of the fat of slaughtered beasts; in the blood of bullocks and of lambs
and of he goats I have no delight. When ye come to
appear before my face, who hath required this at your
hand? Ye tread my courts\dots{} I
care not for iniquity and the solemn day. Your new moons
and your appointed feasts my soul hateth; they are to me
a burden, I am weary of bearing them. When ye spread forth
your hands, I hide mine eyes from you.
Though ye make your prayers many, yet I will
not hear you: your hands are full of blood. ---
Cease to do evil, learn to do good;\dots{} procure
right for the fatherless, plead the widow's cause; come then,
and let us reason together, saith Jehovah,
and though your sins be red as purple,
they shall be white as snow'' (cf.~Isaiah ch.~48;
ch.~58).

It goes in the spiritual world as in the
bodily, that a motion which is to overcome hindrances
presupposes a pressure and a driving force. The terrible
threats in union with the great promises
furnished, as regards the fulfilment of the law, the pressure and the driving force.
Had the promises been exclusively
bound to an outward fulfilment of the outward commandments
according to the letter, then the driving force would soon have been used up,
for experience taught that a just requital
did not take place in the outward; but the prophetic
word held doubt down by placing the guilt
in the disposition, pointing to those who are righteous in truth
(Ezekiel 14:14), and referring to the future. Had
% --- p. 52 ---
\setcounter{footnote}{0}%
\opage{52}the promises been exclusively bound to the inward
fulfilment, then the firm Mosaic foundation would
have dissolved into an indefinite spirituality, a
formless inwardness, which conflicted with the purpose of revelation: to educate the people to a real recognition
of Jehovah's will, to unconditional obedience to the laws
of the spirit.

To understand the old prophets, we must rightly
note this peculiarity, that the fulfilment of the law
points just as much outward as inward, and that the
same holds of threats and promises, while the
outward nevertheless is and remains so heterogeneous with the
inward that the miraculous everywhere shines through.
The psychological tension which thereby arises
and sets the imagination of faith astir casts the light of understanding
on the breakers and surges under which
the prophetic discourses are continually moved between fact
and image, reality and symbolism. --- ``Woe
unto them that rise up early in the morning, because they hanker
after drink. They sit far into the night, and wine
inflames them. Harps, citherns, tabrets, pipes and wine
are at their feasts, but the work of Jehovah they will
not see. Woe unto them that say of evil: it is
good, and of good, it is evil. They make the
ungodly righteous for the sake of a reward, and take away
the right of the righteous from him. Therefore, as the flame
of fire devoureth the stubble, and dry grass sinks together
in the blaze, so shall their root be as mouldiness and
their blossom pass away as the dust; for they despised
the law of Jehovah Sabaoth, they scorned the word of the Holy One
of Israel. Therefore the wrath of Jehovah blazes up against his people;
he stretcheth out his hand over it, he smiteth it
% --- p. 53 ---
\setcounter{footnote}{0}%
\opage{53}so that the mountains tremble: the corpses lie like refuse in
the street.'' (Isa. ch.~5.) This alternating play of image and
fact, symbol and reality is characteristic
of the messianic prophecies.

\emph{Images of the Messiah}. In the law the relation to God is
posited unconditionally as a relation of will; from the
strict demands of the fulfilment of the law spring the image-rich portrayals
of the expected Saviour. At first it looks as
though the people could really fulfil the law, if only it earnestly
would; but gradually it becomes evident
that the same God who has given his people the law
must also by his spirit help it to have the law
fulfilled. The relation between Jehovah's supernatural will
and human freedom is then more closely determined
thus: that the law is transgressed is men's own
guilt, and the divine punishments are just; but that
the law is fulfilled is conditioned by God's grace, God's
spirit, for of himself man has not the strength to
satisfy the law. The actual state is now this,
that the people, on account of the transgressions, finds itself in
distress and danger, oppressed and threatened by overpowering enemies.
The condition of salvation is that the people acknowledge its guilt and
its impotence, turn to Jehovah and trust unconditionally
in his promises. When, then, the prophet portrays
sin and punishment, he keeps close to the actual,
even if the features appear enlarged in the mirror of revelation;
but if he speaks of grace and salvation,
then the imagination glowing in hope goes far out
beyond the natural bounds of reality; the light from the
supernatural breaks forth with many-coloured rays and
glorifies the whole of nature. ``The people that walked in
darkness beholds a great light, a light is risen over them
% --- p. 54 ---
\setcounter{footnote}{0}%
\opage{54}that dwell in the land of the shadows of death\dots{} For unto us a child
is born, unto us a son is given. The government shall
be upon his shoulder; his name shall be called Wonderful,
Counsellor, The mighty God, The everlasting Father, The Prince
of Peace. Of the greatness of the government and of the peace
there shall be no end''. --- The throne of David and his kingdom
shall be established with judgment and with justice (Isa. ch.~9).
If we ask who this Prince of Peace is, and whence he
shall come, the prophet answers (Isa. ch.~11):
``There cometh forth a shoot out of the stump of Jesse, and a sprout from
its roots shall blossom forth. Upon him resteth the spirit
of Jehovah, the spirit of wisdom and understanding, the spirit of counsel and
might, the spirit of knowledge and of the fear of Jehovah\dots{}
He smiteth the earth with the rod of his mouth, and with the breath of his
lips shall he slay the ungodly\dots{} Then
shall a wolf abide with a lamb, and a panther lie
with a kid; a calf and a young lion shall be
together: a child can lead them.''

It is the spiritual Israel from whose root the Wonderful One
springs forth. The fundamental features of his being,
his working and his temporal lot can be seen in the history of the chosen
people; for it is this people, and this
alone, that has had to bear and suffer for the victory of the spirit-will,
for the holy cause of truth and righteousness.
``He grew up as a shoot before the face of Jehovah, and
as a sprout out of a dry ground. He had no
form nor glory, that we should look upon him,
no comeliness, that we could have delight in him. He
was despised and forsaken of men, a man of
sorrows, acquainted with sickness, as one from whom men
hide their face; so despised that we esteemed him as
nothing. Surely he bore our sicknesses, and our
% --- p. 55 ---
\setcounter{footnote}{0}%
\opage{55}sorrows he took upon himself; we esteemed him as one
that was stricken, smitten of God and sorely afflicted. And
he was pierced for our transgressions
and smitten for our sins. The chastisement was upon him, that
we should have peace, and by his wounds we are healed.''
(Isa.~53:2--5). It is the spiritual Israel which,
to the prophet's sight, takes on personal form; it
is ``the servant of Jehovah'', who in patience suffers for the
carnal Israel and takes its transgressions upon himself;
it is he who in voluntary obedience fulfils the law
and brings the great sacrifice which alone has power to blot out
the guilt.

The servant of Jehovah is the son of Jehovah; but the son who
was to restore Israel's kingdom and raise up David's fallen
tabernacle,'' (Amos 9:11) had to have not only a servant's
patience, but also a kingly, indeed supernatural power.
``I will give thee the peoples that rage, and the wide
circle of the earth for thy possession; thou shalt break the heathen with
a rod of iron, like a potter's vessel shalt thou dash them in
pieces.'' (Psalm 2). During the Babylonian exile,
when there was no prospect at all of a restoration
of the kingdom by worldly weapons, faith
in the supra-worldly had to be doubly strengthened. What is
not to be seen on earth the prophet then sees in the clouds
of heaven. ``I saw in the night visions, and behold: with
the clouds of heaven came one like a son of man, and
came to the Ancient of days and was brought before
him. And there was given him dominion and glory and a kingdom, that
all peoples, tribes and tongues should worship him.
His dominion is an everlasting dominion, which
passeth not away, his kingdom shall not perish.''
(Daniel 7:13--14).

% --- p. 56 ---
\setcounter{footnote}{0}%
\opage{56}However unlike these images of the Messiah --- the servant in
lowliness and the mighty king, the meek Prince
of Peace and the supernaturally armed warrior, the
one who suffers on earth and the son of the Ancient in the clouds
--- however unlike these images are among themselves, yet
the family features in them all are so essential that they
fall together into one fundamental image, the image of him who
fulfils the law.

But at the Old Testament standpoint the law
cannot be fulfilled. To fulfil the law is one with
abolishing the law; that is a secret of revelation
which none of the old prophets has been able to see through.
The fulfilment of the law is conditioned by the transformation of Judaism,
a transformation in which it perishes.
As long as the supernatural and the natural, the divine and
the human being, continued to stand
so sharply against each other as they, for the unambiguous
manifestation of the spirit-will, had to stand precisely at the stage of Judaism,
so long must the fulfilment of the law in
spirit and truth remain an unfulfilled demand. After
the time of the exile, when Zion was again seen in Jehovah's light
and the second temple was raised, idolatry ceased,
and the gross apostasy. The law was read aloud, the people
listened and understood that its life and independence unalterably % print: Selstændighed (misprint)
depended on a punctual fulfilment of the law. But
it was really only a fulfilment in the outward, a
ceremonial priestly fulfilment, that was reached:
the spirit was gone, the prophetic voices had fallen silent.
So men lived on upon a tradition.
Pharisaism and Sadduceeism strove over the interpretation of Scripture,
and political ideals from the heroic age of the Maccabees
mingled with the messianic expectations. In order
% --- p. 57 ---
\setcounter{footnote}{0}%
\opage{57}not to be dissolved into Greek and Roman heathenism,
men held, under the pressure of worldly superior power from without,
to the letter of the law, to the forms of worship, and the forms
crystallised. Such was the state of things that
coincided immediately with ``the fulness of time''. The Messiah
came. The spiritual Israel, the narrow circle of
the quiet in the land, greeted him in faith and hope with
the humble joy of expectation; but the priests and
Pharisees were offended at him, feared him, hated
him and saw in him a seducer of the people, a blasphemer,
a criminal against the law. He who fulfilled the law
was judged according to the law and crucified as a malefactor: a
soul-shaking contradiction, deeply grounded in the essence
of Judaism! The law is fulfilled; but the fulfilment is surrounded by
terrors. The Messiah dies: the sun of Judaism is darkened,
its foundations tremble, its graves open, the veil
of the temple is rent; but --- the law was fulfilled,
the prophecies were fulfilled.

% --- p. 58 ---
\setcounter{footnote}{0}%
\lecture{Fourth Lecture.}\opage{58}The reconciliation of the supernatural with the natural,
of the will of the Spirit with the essence of man, was --- on a
supernatural foundation --- accomplished by the God-man.
That God was in Christ, reconciling the world unto himself,
means that he reconciled man's determination as
spirit with his original determination by nature. It is
not a part of man, his heart alone or his
head, his feeling or his understanding, his imagination
or his will, but the whole man, man
as he proceeds from the hand of nature, that is to be transformed.
The whole of nature, and then in a special sense
the nature of man, is the workshop of the spirit, the foundation
on which it builds, the instrument with which it
works, the material it needs, the immediate
stuff it shapes and forms. But, we repeat it,
the reconciliation is grounded in the supernatural. The birth of Christ
is a wonder, his resurrection is a wonder, regeneration,
the new birth of water and of the Spirit, is a wonder which
no scribe Nicodemus comprehends.

But that the wonder in the New Testament points
to a far more inward union of the natural with
the supernatural than in the Old is seen, among other things,
from this, that the prophetic changes its character. The
% --- p. 59 ---
\setcounter{footnote}{0}%
\opage{59}old prophets stood under the Spirit as almighty
will; it came upon them when it would, and withdrew
again when it would. The prophet from Nazareth,
on the other hand, lives in the Spirit; it is ``abiding upon him'',
it is, as the Spirit of personality, one with his holy
will, as the Spirit of love inwardly one with his
pure nature. The ancients spoke in the name of the Spirit: \emph{thus}
\emph{saith Jehovah}; the One sent by the Father speaks in his own
name: \emph{Verily}, \emph{verily I say unto you}.

At the stage of the Law the kingdom of the Spirit is a supramundane
kingdom, which the almighty will --- no doubt for the good of men,
but yet in the proper sense for the sake of its
own dominion --- wills to establish on earth.
From this standpoint Jehovah can be likened to a king
who wages war with earthly potentates, and who, although
he could crush them all with one blow, yet
for the sake of his name and his honor condescends to
make use of human forces. He then recruits
troops, promises the faithful, who will obey his commands,
high wages, and threatens the cowardly and the defiant, who transgress
discipline, with harsh punishment. They are good terms
he offers; the warriors must regard it as a true
good fortune to fight under his banner. By serving such
a king they serve themselves; for he has land
to distribute, and the land he has promised them is fair:
it flows with milk and honey. The faithful servants
can be assured that, when the victory is
won and the land conquered, he will keep his word and let
them live long and happily in the land. But the kingdom,
of course, is and remains his own kingdom: his fighting men
are only servants. It is thus not really
for the sake of the kingdom, but for the high wages and the good
% --- p. 60 ---
\setcounter{footnote}{0}%
\opage{60}terms in the kingdom, that his servants fight. If there were
another king who offered them higher wages and a
richer land, they would naturally prefer him, as is
also seen from the fact that many who think they have not
received their pay and bread in due time become deserters,
to serve another king. Here there is, on the side of
the obedient as on the side of the disobedient, a finite calculation
of loss and gain, an external settlement between
sacrifice and self-interest.

At the stage of the Gospel human nature is essentially
reconciled with the spirit; here another image fits.
The Father of spirits, who wages war with the world-powers
hostile to the spirit, has a Son whom he loves, and
this Son is a man, a real Son of man in
flesh and blood, and the Son of man is the heir of the kingdom.
Now as the Son fights for the spreading of the Father's
kingdom, he fights at the same time for the spreading of
his own kingdom, and gathers fighting men. But what kind of
fighting men are these? Not mercenaries, recruited soldiers
who are compelled to obedience by strict discipline and
fight for external wages. No, the reward is inward,
the obedience voluntary: the aim of the kingdom is the salvation of men.
So he calls them on the terms of salvation and freedom.
``\emph{Come unto me}!'' Let us fight for the kingdom of the Spirit,
for this kingdom of my heavenly Father is also
my kingdom; I, the Son of man, am the heir of the kingdom.
``\emph{Come unto me}!'' You, my brothers and sisters, are
my joint heirs; the kingdom of the Spirit and of truth and of freedom
is also your kingdom, a true and real
kingdom of men. \emph{Come unto me}! Let us fight
for our own kingdom, that it may conquer, and we with it. ---
Here there is no question of obedience from fear or of
% --- p. 61 ---
\setcounter{footnote}{0}%
\opage{61}external wages; self-interest is excluded; life is staked
so that life may be won. The reward is to enter into
the kingdom; the punishment is to exclude oneself from the kingdom.
Thus the aim of the spirit is expressly set as the aim of man.
Man's spiritual development goes as one
with the transfiguration of human nature by virtue of the supernatural,
a transfiguration in soul and body, a victory
over the corruptible, a resurrection to incorruption
and eternal life.

For the prophecies of the spreading, struggles and victories
of this kingdom, the discourses (Matthew ch. 24,
Luke ch. 21) on the destruction of Jerusalem and the end of the world
provide the typical conclusion. Judaism perishes,
the heathen powers tear one another apart, the abyss of the lie
opens and swallows its prey. But when everything that
does not have the sustaining ground of truth topples and
collapses, the revealer of truth stands transfigured
as the Word itself, the Word of life and of personality, of righteousness
and of love, that is to judge
the world. In the Revelation of John this world-upheaval
is further painted in visions and images and
Old Testament symbols. --- Our characterization
of the prophets of revelation is herewith concluded.

With the end of the apostolic age it looks as if
the stream of revelation stopped, and the supernatural
all at once withdrew; for we cannot believe in
the miraculous communication of the Spirit on which the Catholic
Church would ground its infallibility. The Montanist
prophets it rejected early; Bernard of Clairvaux
was a mystical hierarch, Joachim of Flora's ``eternal
Gospel'' found no entry; Birgitta, the Sibyl of the North,
stands as an exception. It is a
% --- p. 62 ---
\setcounter{footnote}{0}%
\opage{62}priests' church, which with its hierarchical dogma of inspiration must
stifle all prophetic inspiration by stifling the
free movement of the spirit. Not only the insolent, defiant
word of unbelief, not only the wild, fantastic
word of the mystic and the enthusiast, but even the truly rebuking, chastening
word that turns against its own corruption, must
it bring to silence: it has burned an Arnold,
it has burned a John Huss, yes, what is more, it
has condemned and burned a Savonarola.

However, the spokesmen of Catholicism could here
meet us with a weighty objection. ``By breaking
with the Catholic Church you Protestants have, without
knowing it, broken with Christianity and denied
revelation, for you have denied the continued activity
of the supernatural. It is this denial which
your new prophets, the prophets of genius, now at last
make manifest. The denial of the spirit with which they surprise you, the unbelief with which they offend you, is
at bottom only a purer, clearer expression of your
own unbelief, your own apostasy, your own denial of the spirit
from the beginning. They are right against you
and will by degrees sweep you all along with them. The flood
of unbelief is swelling, it grows ever broader and
stronger; it has inflow from so many sources, both open
and hidden, that the inundation will soon
rise to a true Deluge.

You would gladly believe that the supernatural once,
in days of old, showed itself and worked on earth;
but in the present it dare not show itself at all, not work at all.
You assume with us that the apostolic age was a time of
miracles, that the outpouring of the Spirit at the feast of Pentecost was a
church-founding miracle; but that the Church which is founded
% --- p. 63 ---
\setcounter{footnote}{0}%
\opage{63}by a miracle must also be sustained and governed by a
miracle, this you will not assume. Sound logic ought surely to
be able to teach you one of two things: either the outpouring of the Spirit
at the feast of Pentecost is a real miracle, and then
our doctrine of the abiding of the Spirit in the Church and its supernatural
presence at the ecumenical councils is
irrefutable; or, if our Church's dogma of
the continuing inspiration is called into doubt, then
the inspiration at the first Pentecost must also be called into
doubt. Either the accounts of the miracles of the apostles are
trustworthy, and then the miracles which in the course of
the ages have taken place through the great witnesses in the bosom of the Church,
through and with its saints, men and women, must also be
trustworthy; or: if the miracles by which the Church
to this very day is glorified by its holy ones are
called into doubt, then the miracles of the apostle Peter must be
called into doubt, the miracles of Christ called into doubt. Christ's
ascension, Christ's resurrection, Christ's birth of Mary,
the pure Virgin, all this you Protestants must, with
your culture and modern consciousness, call into doubt, and
where are you then? You criticized the Church and believed in
Scripture without criticism; you went further: you criticized Scripture
in order to believe in reason; you went further: you criticized
reason, pure reason, in order to believe in?
--- well, that is a question. Your much-praised progress
in free enlightenment is, in brief, this: you began
by pushing the Church aside in order to believe in the Gospel;
you end by pushing the Gospel aside in order to
believe in criticism''.

We have cited these objections, not in order
to defend ourselves against the Catholics, but in order to come to
clarity about the essence of Protestantism and its fundamental
% --- p. 64 ---
\setcounter{footnote}{0}%
\opage{64}position toward culture and science. The Catholics' assertion
that the Protestant Church must necessarily end
in self-dissolution finds strong support among the spokesmen of
the modern consciousness. ``So long as Protestantism
has not yet let go of the last remnant of
the old faith in Christianity as supernatural
revelation, so long will it, as a medieval
relic, continue to be in conflict with reason and culture.''
--- The world-view that corresponds to faith in revelation
is, it is said, \emph{dualistic}; the rational world-view of the present,
on the other hand, is \emph{monistic}.
The watchwords are: \emph{dualism} and \emph{monism}.

So that the educated public may not have
the trouble of finding out the meaning of these obscure
words, a short and comprehensible explanation is given. The word
\emph{dualism} means everything stupid, unclear, illogical
and obscurantist that one can imagine; the word \emph{monism},
which shines like a star on the storm-helmet of the physicist-philosophers,
means the clever, the clear, the logically
coherent, that which alone can satisfy
free thought. After this explanation it is no wonder
that the educated, who want to keep pace with the times,
turn their backs on dualism in order to follow at a brisk trot the monism
shining on the storm-helmet.

``If monism were not the truth, the future would look
black. To derive nature from spirit,
to assert the spirit-will and set freedom highest, is dualistic
and old-fashioned; to derive spirit from nature,
to deify the drive-will and set necessity highest,
is monistic and genuinely modern. If the modern
consciousness were not one with modern nature, if
it were thinkable that the storm-helmets could be blown off
% --- p. 65 ---
\setcounter{footnote}{0}%
\opage{65}the physico-philosophical geniuses, as the downy caps
in the field are blown off ``the dandelions'', what then
would become of freedom and progress? Should the spirit-will
really be higher than the drive-will, and necessity
be mastered by freedom, then the cause of mankind
were lost, old age without comfort, manhood
without courage, youth without hope. If one dares
to doubt monism, then one dares to doubt
Darwin's grand hypothesis, the only scientific
dogma about whose unconditional validity
free thought has not yet dared to doubt. Could
one, were it possible, could one doubt the monistic unity
of human nature with ape nature, the monistic descent
of the human form from the ape form --- dreadful
thought! then one would surely have to despair of the aim of life;
for when man does not know his descent,
by what should he then know his aim; when he
does not know his birth, how should he know
the nobility of his birth? If necessity were not
the blind mother of seeing freedom, if the gigantic exertions
of the human will with virtues and vices,
historical exploits and poetic crimes could not monistically
be traced back to the monistic womb of
monistic nature: what then would become
of the natural rights of man, of the harmony
of life in the endless struggle for existence!'' ---
This is in outline the monistic program of the modern consciousness.

The world-view decried as dualism is,
moreover, not dualism, namely not dualism in the
sense that existence should fall apart from itself into two
pieces, a supernatural and a natural;
% --- p. 66 ---
\setcounter{footnote}{0}%
\opage{66}faith in revelation is, by its principle, not dualistic in the sense
that the system of the laws of nature should fall outside
that of the laws of spirit, and for that reason now and then be arbitrarily broken and
arbitrarily healed by a supernatural will of omnipotence; such
a presentation is a distortion that serves only to
recommend and glorify monism. The
dualism in which the Protestant faith in revelation is
grounded has the unity of spirit and nature, freedom and
necessity, God and world as its presupposition. This
original, all-embracing unity is a unity of
strong oppositions; but precisely thereby a strong and
deep and rich original unity. The unity of
spirit and nature, of God and world, is a unity of
freedom, an almighty unity of will. Hereby is determined
the relation between knowledge and faith.

Protestantism recognizes the right of free thought
to the widest extent, that is, it recognizes everything
that thorough investigation is able to confirm; it
protests against everything that is demonstrably hollow and % print: prostesterer (misprint)
groundless. --- ``Is there a living God who cares like a father
for his creatures, especially for man, or is the
all-producing being, in infinite enjoyment of its
productions, infinitely indifferent to the weal and woe of the living
individuals?'' No human being knows it,
no scientific investigation is able to decide
it. ``Are souls immortal, or do they vanish with
the dissolution of the body like lights that are put out?'' No one
knows it, no scientific proof can decide it.
But when science cannot decide which of the
opposed assumptions is the right one, when the one of
them, according to the verdict of science, is just as possible
as the other, then it is surely folly to demand
% --- p. 67 ---
\setcounter{footnote}{0}%
\opage{67}a scientific decision of that which science cannot
decide. --- But if, then, the educated of the present were in
consequence of this to agree that what cannot
be known cannot be believed either, then Protestantism would
protest against all the educated of the present. To doubt
everything is despair, to live on loose opinions is looseness;
the seriousness of life demands an answer to the questions of life,
a decision on its highest concerns.
``But who is to answer the great questions, who
is to settle the highest concerns?'' Man himself and
no other. ``How many mistakes and errors
are not possible here, then?'' Here only one mistake is
possible, one only, which consists in man's being
mistaken about himself; here only one error is possible, one
only, and it consists in man's going astray in
himself.

By making the highest questions of human life
personal questions and granting man the capacity
and authority to answer these questions himself,
Protestantism has made faith a matter personal in the highest sense
and placed the center of gravity of existence in
personality. But, it will be asked, what has that
to do with positive religion? The answer is: personality
is conditioned by self-understanding, but self-understanding
is conditioned by instruction, upbringing and
means of awakening which positive religion offers.
At first it seems as though faith rested on its authority
alone, but as personal development advances,
the relation is reversed, so that the authority of religion
in the last instance comes to rest
on personal self-understanding. The principle of faith
then coincides with the principle of self-understanding:
% --- p. 68 ---
\setcounter{footnote}{0}%
\opage{68}authority is only a support, a railing on the heights, so that
the one climbing up shall not grow dizzy. Should anyone think
that such a conception of personal faith cannot possibly
be consistent with the faith in a historical revelation
which the Protestant Church has professed and
must still constantly profess in order to be a church:
then let him note well that it is not self-understanding
that produces religion, but the reverse.
There is a world of difference between revealing
the personal truth and recognizing the truth in
the self-revealing personality. The spirit-will is
the center of personality; if the spirit-will is in its essence
only relative, a transient magnitude that arises in
time and passes away in time, then the personal has only
a relative worth, and freedom perishes in necessity.
If freedom is to conquer and personality to count
unconditionally, then the spirit-will must itself be a will that
bears the whole of existence, an almighty will that expresses
itself in the laws of nature as well as in those of the spirit.
Only the man whose personal self-understanding goes
so deep that he recognizes his powerless will as inspirited
and moved by the almighty will reaches the anchor-ground
of personality.

From this it will be understood with what caution the
inner, individual experiences must be treated, how carefully
they must be compared with the outer ones, what
perseverance in psychological observation and self-criticism
is required, in order that the Protestant faith in revelation
can be said to be a fully conscious faith, a faith
that can give an account of its origin and its character.
The Protestant consciousness at its present stage of culture
by no means finds itself in immediate and sympathetic
% --- p. 69 ---
\setcounter{footnote}{0}%
\opage{69}understanding with the spiritual phenomena it meets in
revelation. On the contrary; the more externally historically
these phenomena are taken, the more repellent they are.
Only when the external is led back to its internal
does understanding begin. It is the relation of the spirit-will
to the nature of man that provides the foundation of understanding.

The Protestant consciousness of culture is an ideal
product of manifold transformations, long, preceding
labors; in these labors it has its history;
in this history it must seek the presuppositions for
an understanding of itself. The same holds of the
Protestant consciousness of faith. It has in revealed
religion and the subsequent conceptions of its
character and essence its history, and in this history
the presuppositions for the understanding of itself. Religion
is a relation to God. That this relation to God does not
come about by man's himself seeking and
finding a God, the consciousness of culture can testify.
It is, after all, precisely assured that one can seek God
inwardly and outwardly, on earth and in heaven and wherever
one likes, and yet find him nowhere. Criticism makes it
manifest that the relation to God which is arbitrarily posited by
man himself is subjective and groundless, that the God
who is produced by human feeling, human imagination
and human thought must just as surely become an idol
as the gods that are produced by human hands.

As surely now as the relation to God is the fundamental relation
of personality, and this fundamental relation cannot have
reality unless it is originally posited by God, so
the Protestant personality that is sure of itself
will hold itself assured that the will of God, the almighty
spirit-will, continually stands in an active relation
% --- p. 70 ---
\setcounter{footnote}{0}%
\opage{70}to the human will: this active relation is the wonder.
The wonder, so understood, is not the same as
the miracle. All being is wonder (alles Sein ist Wunder),
this saying of a famous physicist we apply here in the
sense that a divine thinking, a divine
willing pervades the whole material world, that
the forces of matter act precisely as natural causes
in acting according to their natural dispositions, that these natural
dispositions must be referred to a divine
willing, that things, in other words, realize the Creator's
will by acting according to their nature. On the
presupposition that the spirit-will is active in the whole of
nature, it must then also be active in man,
and its activity is seen in the peculiar dispositions of human nature.

From this standpoint the Protestant personality
can, for religious-psychological self-understanding, without
coming into contradiction with science, make clear to himself the
historical presuppositions on which his faith rests.
That there must lie in the nature of man an original
need of religion is seen already from the myths. In the
mythical relation to God the divine is just as
much active as the human. The time is past
when the mythical gods were regarded as arbitrary inventions.
But under the form of myth the divine is
hidden in the natural; only at the stage of revelation
does the spirit separate itself from nature, as
seeing freedom from blind necessity; the relation to God
is posited by the spirit itself as a relation of will,
and the consciousness of the personal comes into being for
man.

% --- p. 71 ---
\setcounter{footnote}{0}%
\opage{71}How it actually came about that the divine
spirit, in the course of the ages, in the progress of development,
changed its willing into will, how
eternal freedom stepped forth against natural necessity and % print: Naturnødvendighen (misprint)
set its self-revelation in the place of myth, these are
questions whose answer lies beyond the bounds of
knowledge. The objection that no change can
take place in the divine mode of operation, since
God surely is and must be unchangeable, is reasoning
in the air; it lies precisely in the nature of will
to be able to change its operation without changing its essence.
The mystery is the unfathomable relation between freedom
and necessity; we must keep to the religious-psychological.


As the Protestant personality feels repelled
by the miracle of the Spirit, he discovers in what is immediately
repellent at the same time something essentially attractive. It
is the speaking God, the supernatural declaration of God's will,
that repels; and why? Because such a
declaration falls outside all experience in
the present and presupposes a state of soul that has long
since vanished. But in this strange thing there is yet for him
something familiar, something inwardly akin to his own state.
At his own stage of consciousness the personality
cannot in and by himself discover the influence of the supernatural
will; in his inner experiences he is referred
to a general presupposing of divine
thinking and divine willing; he does not get, and
cannot get, further, for his relation to God is reflected.
What then does he discover? That this reflected
relation to God is a mirroring of that original one, and
religious inspiriting a psychological after-effect
% --- p. 72 ---
\setcounter{footnote}{0}%
\opage{72}of the miracle of the Spirit. He does not believe in the miracle of the Spirit
because he believes in the authority of a tradition; but
he believes in the authority of tradition because he
recognizes in self-understanding that the perennial wonder
in which his relation to God is grounded, the increase of will
by which his independence grows, the luminous ether
in which his freedom breathes, includes the miracle of the Spirit
as the possibility of its coming-to-be, its original presupposition.
So understood, the derived leads to continued
deepening in the original.

The Protestant faith in revelation is essentially
different from the Catholic. The difference comes
only partly to light so long as the dispute is chiefly
about what it is \emph{in which} one is to believe, \emph{how much} one is to believe, etc. The fundamental difference
shows itself only when the question is \emph{how} one
is to believe, and what faith is in its principle.

Catholic faith is immediate faith in authority,
its faith in miracles a spiritless faith, and spiritless faith
superstition. The Catholic Church has not merely imitated
Judaism in externality; with its enormous
apparatus of tradition, its forms of doctrine, its cult, its
church discipline and above all through its intricate world-politics
it has gone much further. Those who love prophecies
in the form of propositions and let the prophet ``behold in the spirit''
the future when he speaks of his own day, can
rightly assert that it is Catholic externality
that Isaiah beheld in the spirit, since they can here repeat
with emphasis, yes, with tenfold emphasis,
the words of Jehovah: ``I am weary of your many
sacrifices\dots{} Your new moons and your appointed feasts
my soul hateth; they are a burden unto me.

% --- p. 73 ---
\setcounter{footnote}{0}%
\opage{73}At last, when the abuse was at its height, there came
an outbreak from the ground, and the forces of the spirit stirred.
But the breakthrough was no miraculous
outbreak; it was no new revelation, it was a return
to the original, it was the Reformation.
Luther is a prophet, but not by virtue of a miracle of the Spirit.


In Luther's inspiriting the supernatural does not come
specially forward; it is on the contrary hidden under sheer natural conditions.
In judging Luther's grand
character, the point is to conceive him at once
as a man of the spirit and a man of genius. In a
prophet like Isaiah, an evangelist like John, an
apostle like Paul, the greatness of natural endowment is unmistakable:
but one will not on that account conceive Isaiah,
John and Paul as geniuses. And why not?
Because the natural in them is so wholly in the service of the spirit-will
that it does not shine with its own light,
but feeds the flame of the spirit and only gives it a
changing tint. In Luther the relation is reversed;
the spirit-will is here and stirs powerfully. Luther has
courage because he has faith, and he has faith because he
is deeply gripped by the spirit. Read, for example, his letter
to the Elector when he left the Wartburg: ``Your Electoral
Grace knows, or if you do not know it, I hereby make it
known to you, that I have the Gospel not from
men, but from heaven alone through our Lord Jesus
Christ\dots{} But that I offered to submit
myself to examination and judgment happened not because I
doubted, but out of humility, to entice the others.
But since I now see that my all too great humility would
turn to the abasement of the Gospel, and that the devil would
% --- p. 74 ---
\setcounter{footnote}{0}%
\opage{74}take the place altogether if I yield him
but a hand's breadth, I must, compelled by my conscience,
take hold of the matter otherwise. I have shown Your Electoral
Grace obedience enough in that, in order to serve
Your Electoral Grace, I have held back this year\dots{}
This I know with myself, that if this matter stood
in Leipzig as it stands in Wittenberg, I would
nevertheless ride in there, even if (may Your Electoral Grace
forgive me my foolish talk) it rained nothing but
Duke Georges for nine days, and each of them were nine times as
furious as this one\dots{} Let this be written to Your Electoral
Grace to the end that Your Electoral Grace may know
that I come to Wittenberg under a far higher
protection than the Elector's. Nor have I any
mind to ask for Your Electoral Grace's protection; indeed I
hold that I should sooner protect Your Electoral
Grace than you could protect me. Moreover,
if I knew that Your Electoral Grace could and would
protect me, I would not come. This matter
shall not and cannot be helped and furthered by the sword;
God alone must govern here, without any human care
and assistance. Therefore he who is strongest in faith
is here the best protector; and since I now perceive that Your Electoral
Grace is quite weak in faith, I can by no means
hold Your Grace to be the man who should be able to protect
or save me''.

When one considers that Luther, when he wrote
this letter, was an outlaw, and that the Elector Frederick
the Wise was the only one who, humanly speaking, could
take him under his protection, one gets the impression
of a faith and a courage that recall the old
prophets.

% --- p. 75 ---
\setcounter{footnote}{0}%
\opage{75}Nonetheless the activity of the spirit is peculiarly
conditioned by his individual nature. His
crises of soul, the conditions of his age, the temper of his people,
his grand passion, the strong lightning from
the spiritual storm into which he comes, all together
belong to a cultural-historical atmosphere that surges about
his figure and shows him to us as a man of the idea,
a man of genius. However independent Luther was of
all calculations of the understanding and finite considerations of prudence,
he was nonetheless unable, with all his strength of faith and
the honesty of his religious character, to hold fast his
evangelical idea of freedom in its original purity;
he had at last to give way to the conditions of the world and
put up with limits he was not able
to break through.

It is remarkable to see how the priestly-theological
element in Luther little by little gains preponderance
over the prophetic. Zwickau prophets, peasants' revolt,
heretics and enthusiasts of every color showed by
clear signs that the wild spirits were let loose. ``When
we'', says Luther, ``with the help of the Gospel began
to teach that external things could not save,\dots{}
these people, even great and learned people, went so far
that baptism as external water, the Word as external,
human speech, Scripture as external
letters made of ink, wine and the bread as baked by
the baker, were to be nothing at all, for they were external,
perishable things; so they fell to the cry:
Spirit! Spirit! The Spirit must do it, the letter killeth!
--- But ``Paul'' (1 Cor. 15, 3--4) regards
Scripture as his strongest testimony and shows that there
is no means of holding our doctrine and faith except
% --- p. 76 ---
\setcounter{footnote}{0}%
\opage{76}the bodily or written word, composed in letters
and orally preached by him or others.
For here it stands clearly: Scripture, Scripture. But Scripture is
not sheer spirit, about which they babble that the Spirit must
do it alone, that Scripture is a dead letter and cannot
give life. Make no great ado about the Spirit
when you do not have the external, manifest Word, for it
will surely not be a good spirit, but the loathsome Satan
from hell. For the Holy Spirit has surely enclosed
his truth, his counsels and all mysteries in the Word
and revealed them in Scripture, so that no one can excuse
himself or has to seek or search after anything
else''.

However gladly we grant the validity of these
words when they are rightly understood, however much we
appreciate the endeavors of Luther and the Reformers
to secure for the Church a confession of faith and to build
the system of doctrine on a firm foundation, we must yet
expressly emphasize that it was precisely on a mistake
in these endeavors that the carrying through of the principle of Protestantism
ran aground. What lived in
Luther himself, but never dawned on Luther's consciousness,
the great problem of how things really stand
with the Word and the Spirit, with the authority of Scripture and
the free appropriation of truth, could become still less
clear to his zealous followers. They yearned
for a new state of things and worked for something established;
they needed order and discipline in what is external, and out of
the words of Scripture and the teaching of Scripture, whose elements the
heat of the Reformation had melted, and the wild spirits
had mixed with impure ingredients, they crystallized out the
Lutheran dogmas for juridical use and formed a firm
% --- p. 77 ---
\setcounter{footnote}{0}%
\opage{77}foundation for the Lutheran state church. Once more faith in authority,
once more externality!

Protestant externality was in the seventeenth
and eighteenth centuries just as spiritless as the
Catholic, and far more outrageous, because it stood in
glaring contradiction with the principle. As the
glowing masses and hot vapors in the interior of the earth, under
the pressure of the solid crust, increase in elastic force until
the crust cracks and a volcanic eruption occurs, so
the inquiring spirit, under the heavy pressure
of external authority, at last gained the elastic force to
break the crust and give itself air. After single, scattered
eruptions the flames of freedom break out in a chain of
volcanoes, serial volcanoes, whose fire cannot possibly be quenched;
for when one stops up the crater of one, the
others are in full activity: in the glare of these volcanoes
we meet \emph{the new prophets}.

% --- p. 78 ---
\setcounter{footnote}{0}%
\lecture{Fifth Lecture.}\opage{78}The French Revolution of 1789 was not one rebellion
among many; it was a revolution of the foundations,
an earthquake in the history of the world. The tremors reached
far and were felt everywhere, not only in what is outward,
but also, and chiefly, in what is inward. That
rotten pillars fell, that thrones toppled together with old dynasties,
was perceived with terror in the visible world;
but in the invisible world too, the world of minds, of thoughts, of ideas,
there were dynasties that fell: authority fell,
legitimacy fell, piety fell, the venerable
majesty of Christianity was trampled underfoot. It was
the drive-will from the abyss of blind necessity
that raised the banner of freedom and rose up with its
whole army of selfishly flaming passions. A stream of fire
swept through the wild souls, and human heads
were mown down like thistle heads. The demon of the Revolution
shook its bloody wings and swore not to rest until
all the barriers to freedom had been broken through, the bulwarks
of privilege razed, and the towers of tyranny
levelled with the ground.

When the storm had at last spent its fury, when the old
altars with the old faith, the old thrones with the
old authorities had risen again, there were no doubt
% --- p. 79 ---
\setcounter{footnote}{0}%
\opage{79}many who thought that a repetition of such dreadful
catastrophes could be prevented, and the traces of the devastations
wiped out, if the legitimate rulers
took the grace of God to their aid, and Church and police joined
in applying greater vigilance; but experience
has so far not confirmed their expectation. We shall
not treat the political complications here, but fix
our attention exclusively on the ``Kulturkampf'',
the struggle that is waged between the old and the new
with ideal weapons. The old has no living
prophets, all its prophets are dead; but --- the new
prophets are living.

Which, then, are these new prophets? What do they
look like? Where do they come from? What do they want? Where is
the goal to which they point? It is difficult to
answer these questions precisely. The crowd is
motley, its number scarcely to be counted. Different
physiognomies, masks, costumes: how different they are, after all,
in gifts, education, bearing and attitude!
They come, indeed, from all spheres of life, from the workers'
workshops, from the camps of politics, from the hearths of agitation
and from the quiet realm of science; we
find them in literature, we hear them in the popular assemblies;
they shine in practical adventures, we meet them
in novels and plays. A Byron and a Darwin, a
Saint-Simon and a George Sand, an Auguste Comte and a
Hartmann, a Louis Blanc and an Alfred de Musset, a
Proudhon and a Stuart Mill, a Fourier and a Lassalle,
how utterly unlike, indeed how much at odds with one another! Nevertheless
they all belong under the same category; they all stand
on the ground of the Revolution. Despite their mutual
divergences they are agreed in one thing, and this one thing is decisive: the
% --- p. 80 ---
\setcounter{footnote}{0}%
\opage{80}supernatural must go, only the natural counts;
the spirit-will is a shadow, only the drive-will has flesh
and blood.

Without going into the differences in
detail, we shall here confine ourselves to what in
this connection concerns our task, a characterization
of the common main thought, and, for closer illumination,
cite some examples.

\emph{A new religion}. As the first in the series
we must name Saint-Simon, the founder of the religion
called Saint-Simonism, a characteristic product
from the hearth of the Revolution. When religion is spoken of here,
it must be well noted that it is religion in the modern sense,
a religion without God. --- Claude Henri, Count
of Saint-Simon (born 1760, died 1825) belonged to the highest
nobility in France --- he claimed descent from Charlemagne
--- and led an adventurous youth; was Lafayette's
comrade in arms in the American War of Independence, etc.
Alternately rich and poor, industrious and idle, ambitious
and dissolute, he set himself as his goal: to save the
fallen human race and to do better than Christ.
Christ had provided only for the spiritual and overlooked
the fleshly, which was a great one-sidedness; Saint-Simon
would provide for both the spiritual and the bodily
by restoring the flesh to its rights (Nouveau
Christianisme). Instead of an otherworldly, unreal kingdom
of heaven, he would found an earthly, real kingdom of workers,
sanctify industry, work upon the Godhead
by working upon matter, and improve the lot of the poor.
He thus said: ``If this very day death carried off
all the princes of the blood, all the high officers of the Crown,
the ministers of state, presidents, bishops and 10,000
% --- p. 81 ---
\setcounter{footnote}{0}%
\opage{81}of France's greatest landowners, it would cause grief,
because they are persons of high standing; but the state would
not be shaken by a hair's breadth, and the day after these
30,000 pillars would have been replaced, since thousands of
people are able to accomplish what princes
of the blood, ministers, rich men and high prelates have to
do in the world. But if those who died in this way were
the most important master craftsmen, the greatest employers,
the country's chemists, physicists, artists and poets,
then the loss of these perhaps only 3000 persons
could not be made good at all.'' --- Saint-Simon was set by his
disciples high above Christ. He lived for his
idea, but yet did not die like Christ, for his death
was the consequence of a failed suicide. To the disciples
who stood by his bed he spoke these prophetic
words: ``The pear is ripe, you will pick it.''

And they picked ``the pear'', but soon fell out.
It was chiefly marriage and the emancipation of woman
that provoked the quarrel. One of the high priests
in the Saint-Simonian hierarchy, Père Enfentin, had
namely discovered that in the human race there are two
kinds of male individuals. There are, namely, individuals with
feelings so constant that, once they have
fallen in love, they cling with a firm, stubborn persistence
to one and the same person; these are ``Othello natures'',
peculiarly disposed to lifelong marriage
(le mariage définitif). In contrast to them there are
other individualities, lively, mobile natures, who
flit from beauty to beauty --- chivalrous butterflies,
who do not stay hanging on a single flower,
so as not to do the other flowers an injustice! --- these are
``Don Juan natures'', which are fit only for
% --- p. 82 ---
\setcounter{footnote}{0}%
\opage{82}transient marriage (mariage provisoire). Père Enfentin
confessed openly that he belonged to the latter. If this
was now the case with the male individuals,
then there was surely good reason to presume that it must
be similar with the female;
only that these, on account of long oppression,
did not dare to make their innate natural peculiarity known so
straightforwardly. Père Enfentin therefore found that
it was his duty as high priest to investigate
the matter gently and feel his way forward, in order to discover ``the
free woman''.

\emph{New-moral motives}. The old morality
is wont, when the talk is of virtue and duty, to refer
every human being to his own conscience and
remind him to subdue the dangerous passions. But
from the modern standpoint this morality must be rejected.
The drives, the inclinations, the passions are not to be
subdued; on the contrary, they are to be let loose into free and many-sided
interaction. By subjugating the drives, chaining
the passions and blunting the inclinations,
one robs the individual of its innate incentives, the most powerful
motives for moral activity. Man
cannot do anything other than what he has a desire to do, when
it is to be done to good effect. The one among the new prophets
who has declared this with the greatest energy is Charl.
Fourier (died 1837). In his doctrine of the four movements
(le mouvement social, le mouvement animal, le mouvement
organique et le mouvement materiel), whose second
part expressly recommends itself to the pleasure-seeking (les
voluptueux), he sets forth the fundamental thought that, just as
the physical world-order subsists by material
attraction according to Newton's law, so the new
% --- p. 83 ---
\setcounter{footnote}{0}%
\opage{83}social order will come to rest on a mutual
attraction of the individuals' inclinations (attraction
passionnée); every hindrance in this respect must
be especially painful for woman. The universe is a constantly
creating association, in which all suns and planets
are members and co-workers according to their different
properties and inclinations. On the earth the human race
is the creating element. It is the task of society
to restore the harmony among the manifold drives and passions of men
by grouping the individuals for different labours according to their inclinations
and interests.

When one has rightly grasped that the moral motives
are to be sought exclusively in drives and inclinations,
one will be able to understand the expediency
of the provisional marriages and penetrate into the mysteries
of eroticism. In one of her novels George Sand
has the lover say: ``I have never worked upon my imagination
in order to kindle anew or revive a
feeling in my soul that was no longer to be found there; I have
never imposed love on myself as a duty, constancy
as a role. When I have felt love being extinguished
in my soul, I have said so, without being ashamed
of it and without remorse of conscience.'' The mistress in
another novel from the same hand confesses unreservedly:
``Of all these infatuations to which I had given way childishly and blindly,
none, and no liaison, seemed to
me so blameworthy as the one which, in spite of myself,
I tried to make last beyond its time.''

\emph{Denial} of \emph{God} and \emph{Providence}. If dualism
is intolerable because it presupposes not merely a world-nature
but also a God, then for monism God must
% --- p. 84 ---
\setcounter{footnote}{0}%
\opage{84}necessarily be abolished, for the actual nature, after all,
cannot be abolished. But when God drops away,
Providence stands on weak feet. The prophet and
world-saviour who has made himself famous above all by declaring
open war on God and Providence is Pierre
Jos. Proudhon (b. 1809, d. 1865). A true fiery spirit ---
an Elijah among the modern prophets --- Proudhon
raises his face towards heaven, addresses the highest
being and says: ``Now you are dethroned and crushed.
Your name, which for so long has been the last word of the wise man,
a consecration for the judge, the strength of the prince, the
hope of the poor, the refuge of the penitent, this intractable
name will henceforth be devoted to contempt, to cursing,
and be mocked by men.'' What, then, is
God? he asks. ``God is folly and cowardice, God
is hypocrisy and lie, God is tyranny and misery,
God is evil.'' This pathetic expectoration
finds an even and calm expression in the positivists' principle
that the absolute does not exist at all, that everything is relative.

That there can be no Providence, of that
Proudhon is clear in his own mind, not on account
of the natural evils, diseases, the infirmities of old age
and the like, for these are, after all, consequences of the finitude
of individuals; nor on account of the moral
evils, for these can, after all, be ascribed to men's misuse
of free will; but --- on account of the economic
evils. ``Had there been a divine Providence,
it must surely have known from eternity what we
poor mortals only now come to know after 6000 years of painful experience,
that freedom, wealth, knowledge
cannot be furthered and realized in societies like ours
without plunging us into an abyss of misery.
% --- p. 85 ---
\setcounter{footnote}{0}%
\opage{85}Why has Providence not warned us, and led us on better
ways, but left us to our own imperfect logic?
Because there is no Providence.''

Maternal nature, which must now take the place of Providence,
nurses her children at the breast of necessity
and makes experiments with their upbringing,
experiments that are both fine and chemical. ``Every
nation'' --- says Taine --- ``resembles a great experiment
undertaken by nature. Every country is a crucible
into which different substances are thrown in unequal proportions
and under particular conditions. Those substances
are the temperaments and the characters. These conditions
are the climates and the original situation of the social classes.
The mixture ferments according to fixed laws, imperceptibly
through centuries, and produces in the one
case a solid substance, in another a compound that
explodes. It is a pleasure to observe the dull
labour that brings these gigantic masses slowly
and unceasingly to the boil. One comes to know the countless
forces that crush or scatter or draw together
the multitude of living particles that are
subject to them. One feels the regular progress
which, through a series of stages one can count
and of transformations one can foresee, brings them to
the final and predetermined state. One rejoices
sympathetically at the omnipotence of nature, and one smiles
when one sees the eternal chemist, by a slight change
in proportions, in conditions or in substances,
produce revolutions, fashion destinies, call forth
greatness or decline, and determine in advance for each
people which works it shall accomplish, and what misery
it shall suffer. It is a noble spectacle to see the
% --- p. 86 ---
\setcounter{footnote}{0}%
\opage{86}infinite laboratory extended in time as in space,
where so many different vessels, some cold and
filled with the barren ash, others boiling and red
with fruitful flames, reveal the diversity of surging life
and the uniformity of the immortal laws.
Though shut up in a corner of space and of time,
though ephemeral beings whose life is perhaps cut short tomorrow
by the effects of an explosion or by the accident
of a mixture, we can yet discover several
of those laws and form for ourselves a concept of life in the
connected whole.''

\emph{The improvement of the world}. Having freed
himself from the old prejudices and thrown off the heavy
supernaturalist pack, modern
man feels so light, so light; the chest expands, the brow
is raised, the head clear and the muscles ready: now
the work is to be taken in hand. But it is a hard work.
The world is to be improved, society transformed from
the ground up, and the warped human nature as
good as entirely remade. What is it that makes
the provisional marriages not quite work?
The warped human nature. ``Marriage'', says
George Sand through her dear Jacques, ``is now and
always, in my opinion, one of the most odious institutions.
I do not doubt that it will be abolished
if the human race makes any progress towards
justice and reason; a more human and no
less holy bond will then replace it and will be able
to secure the existence of the children without
chaining the freedom of the parents for ever. But the men are
far too coarse, and the women far too cowardly, to demand a % print: alfor (misprint)
nobler law than the iron law that rules them.
% --- p. 87 ---
\setcounter{footnote}{0}%
\opage{87}For beings without conscience and without virtue heavy
chains are fitting. The improvements of which some noble spirits
dream are impossible to realize in this
century; these spirits forget that they are a hundred
years ahead of their contemporaries, and that before one changes
the law, one must change man.'' --- What
is it that makes Fourier's phalanxes not work?
The warped human nature. ``Men
are to be ordered according to their different desires in
circles of 12--1800, which have their own buildings,
phalanstères. All circles across the earth are to stand
in connection with one another, and the threads of connection
run together in Constantinople. Each can choose whatever
occupation he likes, and when he grows tired of
it, he can exchange it for another. No desire
is to be set a limit, no self-esteem
is to be humbled, but each is to have his share of the yield
in proportion to capital, labour and ability.'' Can
anything better be offered? Lifeless nature is normal,
it still goes, as from eternity, according to Newton's \emph{attraction matérielle}, but so warped is human
nature that it no longer obeys its own
law of nature: it will not, after all, go according to Fourier's attraction
passionée.

But poor human nature is lawfully excused.
What use is it that it has a will,
when the will is a product of opposing inclinations
(egoism and altruism), just as water is a
product of oxygen and hydrogen? The great prophets of the physics-philosophers,
Auguste Comte and Stuart Mill, have after all thoroughly
proved that in ``social statics'' and ``social dynamics''
they know their way neither with
% --- p. 88 ---
\setcounter{footnote}{0}%
\opage{88}natural necessity nor with human freedom, when it comes
to the point. That ``human actions are the
result produced by the general relations of human nature
together with man's own particular character''
is incontestable. But when we then ask how
the particular character forms itself, we get the
answer ``that the character in turn is a consequence of the natural
circumstances that bring it up'', and it is added,
as something that goes without saying, that the individual's
``own conscious endeavours also belong among the circumstances''.
Then one looks at the statistical tables,
and sees how what is peculiar in all
productions rests on the circumstances. The moisture
of the Dutch soil produces the grass,
the grass produces the cows, the cows produce the butter
and the cheese, the butter and the cheese produce the Dutchman ---
on account of the circumstances.

How, now, is one to get such a world improved?
If one lays stress on the self-conscious exertion of the human
will, only that this will
must not imagine that it can raise itself above
natural causes or free itself from the influence of conditions and
circumstances, then one does after all presuppose
something peculiar which one calls will. The improvement of the world
is to take place through the united exertions of millions of wills.
The will must exert itself: nature does not
exert itself. The riddle of the world, then, lies enclosed
in the one question: what is the will? If one answers:
the will is a composite quantity, a product of elements
such as ``virtue and vice, such as vitriol and sugar'', then
we must ask: which of the many elements is it, then,
that must exert itself? Is it only a single
% --- p. 89 ---
\setcounter{footnote}{0}%
\opage{89}element, or are all the elements perhaps to exert
themselves? Are they to exert themselves before the product, with the product
or after the product? --- One will surely be
forced to give up the improvement of the world and let it
rest at rejoicing sympathetically over ``the omnipotence
of nature'', and smiling when one sees ``the eternal
chemist'', by analyses and compositions of the elements,
``fashion destinies, call forth greatness or
decline'' and determine in advance for each people'', and then also
for each individual in each people, ``what great things it shall accomplish
and what misery it shall suffer''.

\emph{Hope} and \emph{despondency}. That the new prophets go forward
with a radiant hope is understandable; they want, after all, freedom
and progress, and freedom and progress is the watchword of mankind.
But that it can nonetheless grow dark for them,
and clouds of despondency spread over their faces, when they
rightly consider how warped human nature is, and
how slowly the improvement of the world proceeds, is
also understandable. It is all very well with all this
sympathy with fate, this reconciliation with necessity,
this resting at the breast of nature; but such things also have
their questionable side. ``The eternal chemist'' with his
experiments is not always to be relied on; nature's motherly breast
is now too warm and now too cold. ``Before the
great and fruitful mother, whom nothing disturbs, because
nothing is outside her, there develops'' --- says Taine ---
``in the soul a feeling, a deep and hitherto unknown one. Its very
foundation comes to light; the countless layers
that life has deposited over it, all remnants of passions
and hopes, all the human filth
that has accumulated on the surface, vanish; it
becomes simple, it finds again the instinct of the most ancient times,
% --- p. 90 ---
\setcounter{footnote}{0}%
\opage{90}the indeterminate, uniform words that once put it in
connection with the gods, the natural gods who live
in things; it feels that all the words which it has since
said or heard are only involved chatter, only outbursts of the
moment, only noise in the street, and that if there
is one minute in life which is healthy and desirable,
it is the one in which the soul forgets the buzzing and vexations
of the great anthill, and, as the ancient sages
said, perceives the harmony of the spheres, the heartbeat
of the eternal universe.''

But nature is not always in so festive a mood, and
the soul not always prepared for such a festival of eternity.
When the fundamental forces break loose and the elements
hammer, the individual must creep into hiding if it does not want
to be crushed. Even in the sublime calm with which silent
nature looks down on the bustle of the anthill, with great and
small, there is an indifference which recalls the
superior smile with which the gods of Olympus once looked
down on the mortal races. ``What a collection of
ruins, what a graveyard is history! How
many human hearts have beaten of which
no other trace has remained than a form
hewn in a stone! How indifferently the
peaceful sky smiles, how cruel is the beauty of the shining
dome that day after day stands spread out
over races that fall, as a canopy stands
over a coffin.''

In softer natures, in an Alfred de Musset,
such a mood of despondency can, in critical, painful
moments, even turn into the blackest despair.
Under the mask of \emph{Rolla} Musset has expressed despair
over the fact that the old has fallen and that the new
% --- p. 91 ---
\setcounter{footnote}{0}%
\opage{91}has no foundation. ``I do not believe, Christ, in thy
holy word; too late have I come into a world
that has outlived itself. From a century without hope arises
another without fear; the comets of our century
have depopulated the heavens. --- Yet let it be
permitted to the most unbelieving child of this century without
faith to kiss the dust upon thy cross and to weep, O Christ,
upon this cold earth, which lived by thy death and
which will die without thee. For whom do you labour, foolish
destroyers, when you mutilate Christ on his \emph{altar}\dots{}
You would shape man after your imagination, you \emph{would}
form a world. Well then, you have formed it. Your
world is splendid and your man is perfect,
the mountains are levelled, the plain is lit up, you have prudently
pruned the tree of life. All is well swept on your railways,
all is great, all is good, but one suffocates in your
air. You make splendid words ring out; they flutter far away
in the pestilent winds. --- Hypocrisy is dead, one no longer
believes in the priests; virtue dies, one no longer believes
in God. Voltaire can then sleep content, with the hideous
smile fluttering over his dried-up cheeks. His
century, they say, was too young to be able to read
him; ours must please him, and his men are born;
for the monstrous building has fallen upon us, the one he with
his broad hands undermined night and day. Well then,
since it is so: long live love, which intoxication
accompanies; let the burning kisses glow with the
wine of Spain. Let the spirit of giddiness and the noisy
feasts lead us into the arms of the genius of joy. Well then,
sing of Bacchus, of love and of madness, drink
to time, which vanishes, to death, to life. Forget,
% --- p. 92 ---
\setcounter{footnote}{0}%
\opage{92}drink, let freedom live, sing of gold and night,
wine and beauty!''

\emph{Nihilism}. One sees from this what the world is, what
life is worth and what all human thinking and striving
amounts to; it is disappointment, delusion and illusion, all of it.
In the author of ``The Philosophy of the Unconscious'',
E. v. Hartmann, the modern world-view,
\emph{monism}, has a prophet who is all the more reliable
as he is at the same time a monist philosopher. There
are, according to v. Hartmann, three stages of consciousness, corresponding
to three stages on the way of the illusions. --- The first stage
of illusion is marked by the individual's imagining
that it can, in the given state of the world, attain bliss
on earth. That this is a deception is easy to
see. All pleasure is only a condition for displeasure, all enjoyment
for suffering, all joy a condition for sorrow and
pain, and conversely. If men would only
make it clear to themselves that joy is related to
sorrow overcome as the cessation of pain to the torments of the rack,
they would soon be cured of the
spreading craving for pleasure. What is nowadays called
the spectre of pauperism is nothing other than the consciousness
that the struggle with want, the sorrows and their relief
lie on the pain side of the zero point of sensation,
while earlier, when the poverty of the masses was ten times greater,
this consciousness was lacking, and the people bore its poverty
as by the grace of God. --- At the second stage of illusion
the individual imagines that it is to reach happiness in a life
beyond, after death. That is the Christian deception.
Jesus's world-view is too naive and childlike to separate
soul and body, which is why ``the resurrection of the flesh''
is also included in the third article of faith. --- The
% --- p. 93 ---
\setcounter{footnote}{0}%
\opage{93}third and last stage of illusion is the one on which we ride as
on a railway. The imagination is that happiness lies in a
great political-social future. One expects salvation
from an improved social order, from progress in industry
and science: a monstrous deception!

``The end will be that the old, outlived human race
is tired, tired to death, tired of the labours and
of the enjoyments, without heirs and without love, tired
of everything. Then it will come, that sublime melancholy
which one usually finds in geniuses and spiritually eminent
old men; then it will come, the pure clarity of heaven,
that cool divine calm which wafts in
Spinoza's Ethics, where the passions are swallowed by the abyss
of reason, because they are clearly and distinctly grasped in
ideas. Then illusion is dead; then hope is burnt out,
for what should one still hope for? Then
exhausted mankind recognizes only one thing that is blessed ---
rest from sorrow and want, painlessness in the pure nothing.''

But may we now assert that the new prophets are all
without exception nihilists? It depends on
what one understands by nihilism. According to the usual
opinion, nothing is the pure zero, the thinnest and
emptiest and most tedious of all: something is always
better than nothing at all. From this point of view the nihilists
would accordingly be people who wanted nothing, did not
mean anything seriously and had no goal
for the attainment of which they found it worth working. That
is a great misunderstanding. The nihilists are people with
a sense for the tasks of actual life, indeed with an almost
greedy appetite for actuality. The Nothing that
furnishes the watchword means only that God is nothing, that the
supernatural, the beyond, is nothing, but not that
% --- p. 94 ---
\setcounter{footnote}{0}%
\opage{94}nature with its forces and laws, not that the present
life with its labours, means and ends is nothing.
Is life then really better than nothing? It depends
on the circumstances. As long as the vital forces offer
themselves, one ought to use them. Ought? Of course,
it is no unconditional duty; everything is relative, the unconditioned,
the absolute is nothing. Each must use his
forces in accordance with his individual
nature, his drives and inclinations, and be prepared
for the consequences. That is why we find the idle and the industrious,
the frivolous and the serious, the cowardly and the brave, the effeminate
and the hardy, the pleasure-seeking and the energetic,
egoists who care only for themselves, and altruists who
sacrifice themselves and work for the whole of mankind, in motley
mixtures gathered under the flag of nihilism. They all
grant that life can be better than nothing, but
they all admit too that it can be worse than
nothing: only the philosophical heads --- but they are,
after all, almost all philosophical heads --- know within
themselves that the whole is at bottom nothing.

This ``at bottom nothing'' is the secret, the
secret that does away with dualism. The flaw of dualism
is a twoness, and the twoness splits
again into multiplicity. Actual nature is infected
with dualism: up stands against down, right
against left, the positive pole of the magnet against the negative,
the acid against the base, the male flower against the female flower;
not even the most fleeting provisional marriage can be consummated
without dualism. Monism wants unity; but
the undifferentiated unity of all is nothing. As long as
something distinguishes itself from something else, as long as one leaf in
the forest is unlike the other, one wave in the sea
% --- p. 95 ---
\setcounter{footnote}{0}%
\opage{95}opposed to the other, so long one has only plurality, not
unity. To reach complete unity we must
think of all inequality, all difference as erased.
When all difference falls away, the undifferentiated
unity remains as the ground; but the undifferentiated
unity is at bottom nothing: consistent monism
is nihilism. If one thinks that philosophical nihilism
is impractical, but that monism is nevertheless the only
world-view which, practically understood, suits the present age,
the only one capable of awakening the courage of youth
and pointing to the goal, then we ask: what sort of
liberator is it, then, who stands there --- over by the goal?
Death, the great nihilist! Death is practical, and
death is nothing. Death levels everything, it satiates ethics,
it digests criticism, it abolishes the principle
of contradiction. In nihilism, thoroughly understood, there is
no difference between good and evil, between truth
and lie. The old prophets of life are dead, the new
prophets of death live --- in order to die; so at last
the new meet the old in death's pure nothing, the infinite
all-unity: that is monistic reconciliation.

% --- p. 96 ---
\setcounter{footnote}{0}%
\lecture{Sixth Lecture.}\opage{96}Despite the strong ferment of the age, there is still
without doubt a majority of fellow citizens who, although they
have a good deal to object to in the established order, turn away
with aversion from the new prophets; the
sober judgement runs thus: When we hear these
men's often wild and offensive outbursts of blasphemy,
when we see the arrogant self-confidence
with which they appear, the recklessness with which they
confuse the concepts of duty and conscience and tear apart
moral bonds, the violence with which they attack
the foundation pillars of society, without yet, by their loose,
mutually contradictory projects, putting anything better
in their place, then we must, with all respect for the talent
and well-meant zeal of individuals, hold ourselves convinced that
they are on a wrong path. Indeed, there are surely those whose verdict of rejection
will find a harsher expression. And yet
we may by no means suppose, as regards the matter as a whole,
that the general public is clear with itself. Tender spots
are being touched, something is being spoken of that arouses
attention; the new prophets have words that ring
in the ears of the present day.

It is quite true that spokesmen who know how to
incite the passions and dazzle the crowd with promising
% --- p. 97 ---
\setcounter{footnote}{0}%
\opage{97}prospects must have a great advantage over those who
feel obliged to inculcate self-limitation
and sober progress. Those who say to the common
man: The road is paved, the time has come, strike!
no doubt win more followers than those who early and
late must remind him to take small steps, since
the obstacles are manifold and the rutted roads hardly
passable. How should the leaders who demand
education, discipline and order be able to keep step with
those whose battle cry is: Solidarity, boldness, swift
decision! Therefore we must by no means imagine that
the storm will gradually abate and the ferment settle,
if one only goes on quietly with the old
and carefully avoids touching the real fundamental questions.
``Just hold back, hold back!'' This
``hold back'' acts narcotically on the nerves of society;
the established order keeps silent about its vital questions until they die, as a
sluggish mother overlies her children until they die. When one is then
once roused, as by a pistol shot, from one's spiritual
torpor, one carries on as if it were a matter of shooting
sparrows with cannons. Orthodoxy flies into a rage,
the great bell is rung: ``Antichrist
is come, he works signs and seduces the peoples; Gog
and Magog stand at the door. Out with the sacraments,
away with reason; the priest alone has authority, the
priest alone dares look the living Satan steadily in his treacherous
eye!''

Above all we must not forget that the gifted powers
also derive from spirit, and that the spirit does not let itself be
bound in outward chains; the heavier the chains weigh, the
more active are these spirits, the more zealously do they
work. We must not forget that on the road of spiritual
% --- p. 98 ---
\setcounter{footnote}{0}%
\opage{98}cultivation there stand questions in plenty, decisive questions,
which the prophets of free thought have raised, but
which the defenders of the established order so gladly go
around? Indeed, there are many who think that this is for good
reasons. We must not forget that in the camps of naturalism
there are men who have a sharp eye for states of the soul,
and speak the clear, penetrating language of thorough
inquiry. When such an investigator gathers
himself for a word of admonition, he may have a good deal to
urge upon his contemporaries; let us hear him.

``There is, as is well known, a considerable difference
between juridical and psychological truth, between
juridical and psychological right. When someone assures us that
he believes every word in the symbolic books of the Lutheran Church,
then he has, if nothing else can be adduced
from his other statements that conflicts with this,
undeniably a juridical right to call himself a Lutheran. But
whether he also has a psychological right is another
question. Suppose that he, following the stream of culture,
has gradually come to live with wholly different
notions of heaven and earth, of the natural and the
supernatural, of the being of God and of man, indeed with
notions that the symbolic books expressly
condemn, then he finds himself, even if he has not
become clearly conscious of the deviations, in a state of
psychological untruth; and if, fully conscious
of the deviations, he appeals to his juridical
right, then we tell him to his face that he lives in a
psychological lie.''

``You Lutherans of the present day, with your positive revelation,
your Bible and symbolic books, do you really live
with Luther's faith, in Luther's spirit, after Luther's
% --- p. 99 ---
\setcounter{footnote}{0}%
\opage{99}direction? If you have Luther's faith in Christ, then
you must surely also have Luther's faith in the Devil, for
that is fully warranted in Scripture. That the old
Lutherans believed in the supernatural in dead earnest,
they proved by their living faith in demonic possession.
If you still believe in demonic possession, then
you must not only find it natural that the modern
prophets laugh at you, but, what is more, you must
break with the whole of more recent culture, the present
reality, in which you nonetheless live and breathe. You stammer and
fumble and seek evasions when one presses you
close about the Devil, and do not notice that in so doing you
stammer and fumble in your faith in Christ.
For Christ's personality receives its significance precisely through
its express opposition to the personal Devil;
for this purpose, according to the word of Scripture, the Son of God was manifested,
that he might destroy the works of the Devil. You speak with
unction of sin and grace, but whether you also, with
Luther and Scripture itself, know what grace works,
how it works, and what sin is, of that it is
permitted to doubt; certain it is that many of the liberties
you now regard as permissible the old Lutherans
stamped as mortal sins. You stand with one foot in
the present, with the other in the past; you have two rooms in
your consciousness, you go from present to past as easily as
one goes in a house from one room into the
other, when there is only a door between them, and the door
stands ajar''.

``You scorn monism because it remains
standing at the natural ground of our life without bringing in
something supernatural; you are edified by dualism, that Christ
fights against Satan, so that the new man, especially on
% --- p. 100 ---
\setcounter{footnote}{0}%
\opage{100}Sundays, is busy strangling the old. But
you would not have the old Adam quite strangled, for you are
humane. Your fathers were hard-hearted, they whipped
the old one and could without mercy see the loathsome
sinner hanged. But you are not hard-hearted, you have
pity on the old Adam, after all he is still so
young! To be sure, the rogue must be punished; but
when it comes to the execution, when he kicks on
the rope and gasps for air, then you find the punishment barbaric.
And the old Eve --- still so young and fair!
how much wrong, indeed how many injuries, has she not
to avenge. But when she sees her Adam's repentant mien,
his pleading look, and hears his sigh, then she forgets
everything, then she forgives everything. She bends over the
humble one, she pats his cheek, she kisses his
mouth, she fans him with the fan of poetry, she cools
his brow with the eau de Cologne bottle of the erotic. See,
such is your Lutheran dualism!''

``It must surely be dreadful to live with
a split being, a soul whose one, sinful part
rightly belongs to the Devil, whose other and regenerate
part belongs to Christ. You go about with a two-part
soul and believe in an eternal life. The Devil with the
old Adam pulls to one side, Christ with
the Holy Ghost to the other. If Christ wins,
then you will be eternally blessed, but --- do not forget it! ---
if the Devil wins, then you will be eternally damned. How
can you live on such terms? I know well enough
that there are those among you who imagine that here
there is no danger at all, since you have baptism and faith.
But the faith that is to bring unity into such
% --- p. 101 ---
\setcounter{footnote}{0}%
\opage{101}disruption, peace and reconciliation despite such a tension,
must have more than Herculean powers; it must be
a miracle-working faith, a faith that can move mountains.
What miracles, then, have you worked, how many mountains has
your faith moved? The real mountains you nicely let
stand; so they must presumably be mountains of fantasy. You have
Lutheranism in feelings, emotions and consoling fantasies.
You fight bravely with the old Adam and win the
one fantasy-victory more splendid than the other.
If the Devil now and then comes up behind you with
a little politics, what does that matter? You have, after all, the rock
in the word of faith, and in baptism a children's powder of grace.''

``You have perhaps read a tale by Hoffmann
called ``Elixire des Teufels'', or, if you have not
read it, then by all means read it, so that you may properly
get an impression of what it means to have one's I, one's
being, dualistically split in imagination. --- A young, gifted
monk named Medardus has found in a monastery a
bottle of the Devil's elixir; he drinks from the bottle, and
at once his being is split, so that he can appear
in more than one form. Entangled in love adventures
and dreadful crimes, it becomes quite
impossible for him to find himself again: ``I am, he says,
what I seem, and do not seem what I am;
an inscrutable riddle to myself, I lie in strife
with my own I.'' After having been in prison and having
been interrogated and acquitted of a murder which he nevertheless
really committed, he is released, and now comes his
life's happiest day, it is his wedding day. He
relates: ``There arose a dull noise in the street, hollow
voices cried out among one another, and the booming
% --- p. 102 ---
\setcounter{footnote}{0}%
\opage{102}rattle of a heavy, slowly rolling wagon could be
heard. I hurried to the window. There, just in front of
the palace, stood the wagon drawn by the knacker's men, on
which the monk sat backwards; before him a Capuchin,
who prayed loudly and eagerly together with him. He was
disfigured by the pallor of the fear of death and the bristling
beard; yet the features of the dreadful double were
only all too recognizable to me. Just as the wagon, which had for a
moment been halted by the crowding
mass of people, rolled on again, he cast the fixed,
terrible gaze of his glittering eyes up at me, and
laughed and howled: Bridegroom, bridegroom --- come up on
the roof --- then we shall wrestle with each other, and he
who pushes the other down is king and may drink
blood. I cried: Dreadful man, what do you want
of me? --- Aurelia embraced me, tore me from the window
and cried: For the sake of God and the Holy Virgin\dots{}
They are leading Medardus, my brother's murderer, to death\dots{}
Leonard --- Leonard! Then the spirits of hell within
me awoke and rose up with the power that
had been given them over me, guilty, deranged sinner.
--- I seized Aurelia in fury, so that she started:
Ha, ha, ha --- mad woman --- I --- I, your
lover, your bridegroom, am Medardus, am your brother's
murderer. You, the monk's bride, would howl perdition down
upon your bridegroom? Ho, ho, ho --- I am king ---
I drink your blood.'' --- He rushes down into the street, frees
the monk, hews and stabs and strikes about him, and
flees into the forest. ``Only the thought of fleeing like a
hunted beast stood fast in my soul. I got up, but
I had scarcely gone a few steps before a man
% --- p. 103 ---
\setcounter{footnote}{0}%
\opage{103}sprang up onto my back and flung his arms around
my neck. In vain I tried to shake him off, I
threw myself down, I pressed myself against the trees.
All in vain; the man tittered and laughed scornfully. Then
the moon broke through the black fir trees, shining brightly,
and the monk's deathly pale, ghastly face, the face of the
supposed Medardus, the double, stared
at me with a ghastly look, as if up from
the wagon. Hee-hee-hee --- little brother --- you must carry me,
now I am always with you, I cannot run, I come
straight from the gallows, they meant to break me on the wheel.''

``Lutherans, do you understand this story? --- Luther
understood it when he sat alone in the night and threw the inkpot
against the wall at the Wartburg. It is a little
story about dualism.''

Having heard how bitterly and sharply the
spokesman of modern monism expresses himself against dualism,
we have come to a turning point where the real
question in dispute: spirit-will or drive-will? can
come clearly forward, and the great matter outstanding show
itself in full light. Dualism is no arbitrary
invention, a whim that one can kill with a
stroke of the pen. Dualism has the testimony of life and of reality
on its side: Old and new prophets, what a
dualism! So long as all spirit has not yet become
nature, so long as elementary nature has not yet
been volatilized into spirit, so long as there is still a
remnant of opposition between spirit-will and drive-will
left, so long does dualism remain in force. The animal
with its drive-will does not notice dualism; but
man with his spirit-will gets to feel it
% --- p. 104 ---
\setcounter{footnote}{0}%
\opage{104}thoroughly. Man is dualistic, the ape is monistic;
should the radical means of abolishing dualism
perhaps consist in reducing man to an ape?

The only serious attempt that the new age has seen
at levelling out the dualistic is evidently rationalism,
not in its older, sober form of the understanding, but
in a more many-sided, more spirited form, familiar with ideas,
as \emph{neo-rationalism}. This endeavour has found
able leaders, gifted spokesmen in Germany, in
Switzerland, in Holland, in Sweden. What, then, does neo-rationalism want?
``With Schleiermacher it will seek the innermost essence
of religion in \emph{the depths of the mind}, in the ground of man's life,
in the proper centre of the spirit, but also bring
the unfolding of cognition and of will into inward connection
with this central relation. From an ethical
standpoint it will bring about a purification and renewal of
a large part of the dogmas belonging to the anthropological-soteriological
circle, of the doctrine of free will, of
sin, of grace and of substitution, etc. It
will seek to attain a deeper connection between divine
dependence and human freedom, to
arrive at a truer understanding of the interaction
between the divine and the human factor in
the process of salvation. It will not merely purify and make inward
dogmatics, transform its outward juridical and
magical notions into moral-religious ones; it will
also deepen morality, part it from its vain
self-righteousness, its superficial Pelagianism, in
short, from its finitude, by sinking the roots of morality
in the depths of religion, man's finite
will in the infinity of the divine will, his
freedom in boundness by the divine. It will
% --- p. 105 ---
\setcounter{footnote}{0}%
\opage{105}thereby abolish the fateful and pernicious
separation of the religious and the moral.''\footnote{Den nyeste Theologies Historie i kort Fremstilling [The Most Recent History of Theology, in Brief Presentation], by Dr.\ Karl Schwartz. Danish translation. Copenhagen 1869. pp.\ 466---67.}

It is in truth a noble endeavour; but the question
is whether it can be carried through. That rationalism
takes no account of the welfare of the old Church is obvious.
It will assert the independence of scientific inquiry
from the authority of Scripture and of the symbols,
and in this we must grant it is right; but when it then
demands that faith in revelation shall submit to
the authority of science, then it undermines, surely, both
the Church and the life of the congregation. It admits no
difference in principle between faith and knowledge. If it is, e.g.,
settled that the historical accounts of Christ's resurrection
have no probative force in science, then
they may have no validity in the Church either. If the
Lutheran dogma of the Lord's Supper, scientifically regarded, cannot stand
before the judgement of criticism, then at once the Church must give up its dogma
and make the use of the sacrament dependent on the verdict
of science. Is such a demand well founded?
Yes, in a radical misunderstanding of the essence of faith.

What is it, then, that the religious consciousness seeks in
the historical accounts of Christ's resurrection? A
testimony that the apostles really believed in
this resurrection, that the Church was founded, and the congregation
baptized, on this faith. Let criticism prove that the accounts,
as we have them, do not derive directly
from eyewitnesses, that discrepancies in particulars
suggest that they passed from mouth to mouth before they
% --- p. 106 ---
\setcounter{footnote}{0}%
\opage{106}were written down: it means nothing; the Church can grant
it, the congregation come to know of it, without being in the least
affected by it. The testimony to Christ's resurrection
is a testimony of believers about faith to
believers; the circle of the testimony is so closed that no
criticism from outside is able to break through it.

Only one objection would have validity, but it
is not taken from biblical criticism, it is taken from physics,
or rather from the philosophy of physics, and runs thus:
``Christ is not risen, for the resurrection is a
miracle, and the miracle is an absurdity, a contradiction,
an impossibility''. If rationalism will now bow to
the authority of this knowledge, in order to satisfy the modern
consciousness, then it must take the consequences.

What science is to acknowledge must be capable of being
proved in science. Rationalism speaks of the
living God, a divine love, a living
providence, a holy, divine, righteous will, and
for the sake of religion and morality will not give up
the immortality of the soul; but, we ask, can such things
be grounded scientifically? Modern knowledge can
just as little come to terms with the immortality of the soul
as with Christ's resurrection. If the love of God is to be
filtered through science, then it will in truth be a
cold love.

Rationalism will free the religious consciousness
from immediate dependence on the forms of doctrine handed down
in the Church. The current mechanical doctrine of inspiration
is designated as ``an intolerably crude husk,
in which the kernel is: the unique value of the biblical books''.
The ecclesiastical dogma of inspiration, it is said, has ``long
been recognized as something that caused embarrassment,
% --- p. 107 ---
\setcounter{footnote}{0}%
\opage{107}as a hampering evil. It is nothing other than an
exaggeration and undue intensification of the truth,
serves no real but only an imagined need
of piety, namely this: to have an absolutely finished
expression of the truth.''

We grant it. Personal liberation from outward
authority is of just as thoroughgoing significance
for religion as for science. It is
rationalism's immortal merit to have fought
for this liberation in the Church; but all the more regrettable
is it that it has misunderstood the essence of liberation.
To want, on the authority of Scripture and the Church, to force the
religious consciousness of the present into a literal acceptance
of the biblical notions of the supernatural
kingdom, and still more, into a literal adherence
to the supranatural mechanism of orthodox dogmatics, worked out down to the smallest details,
is certainly so enormous a blunder that one will hardly be able to
say whether it is the men of free thought or
the men of bound thought who contribute most to undermining
religion. But because the people of Israel once
lived in a wholly different circle of notions from ours,
does it follow from this that it had its spiritual life in sheer
imaginings? Because our orthodox fathers formed
concepts of the intervention of the supernatural in the
natural whose objective validity we can by no means
acknowledge, does it follow from this that the supernatural does not
exist at all, that nature itself cannot have its ground in something
supernatural? If the rationalists' ``living God''
is to mean the same as living nature, then
one does not see what difference there is between neo-rationalism
and naturalism. But if the living God as
% --- p. 108 ---
\setcounter{footnote}{0}%
\opage{108}spirit is to be different from nature, then he must
surely be either a supernatural or an unnatural being.

Neo-rationalism has an essential advantage in that,
with a view to the union of the religious and the
moral, it asserts the humane and chooses its centre
in man's conscience. Starting from conscience, the
actively present inward, it strives ``to draw
religion down from its abstract, supramundane, supramoral
height into the very midst of moral interests and reconcile
it with the whole cultural life of the present.'' Will it succeed?
``The congregations'', says Karl Schwarz, ``have awoken from
their slumber; the unnaturalness, intolerance
and lust for power of the Lutheran clergy have risen to such a degree that they
have called forth the strongest counteraction from the people.
The deeper religious need of the educated demands,
imperiously and fully consciously, a different satisfaction from
the one previously offered. We are still in the moment of expectation
and of hope. But the new age
is approaching, we feel its breath of spirit stirring. What
everything points towards, and what must be striven for by all
clear spirits, is the liberation of the Church from the firm embrace
of the existing system of the state, so that it shall not
be dragged along into all its ill-fated crises.
This liberation is synonymous with a building up from
the ground of the people's consciousness, a free shaping of the
genuine, simple, inward Christianity as it lives in
the people. Only in an organic association proceeding from
the midst of the congregations will all the oppressed and
hidden powers of spirit be set free, the whole culture of the present
be led back into the dried-up ecclesiastical channels, and
the Church itself become the
% --- p. 109 ---
\setcounter{footnote}{0}%
\opage{109}\emph{home} of \emph{truth} and of \emph{life}, proclaimer of the most inward faith of conscience.''
--- Will this prophecy be fulfilled?

If a human being is to find in his conscience testimony
to an eternal righteousness, a holy, divine
love, then he must have a conscience of faith;
for in a conscience of knowledge there is nothing of that. He
who conscientiously wants to remove everything that does not
have confirmation in science, and to free his soul from
expectations whose fulfilment cannot be brought
into full harmony with earthly conditions and unchangeable
laws of nature, will find in his conscience no
trace either of supramundane righteousness or of
divine love. Should he once, in an inspired
mood, let himself be carried away into reaching for the
supramundane, he may well be ashamed of himself
afterwards and feel reproaches of conscience; --- ``it
was only a mood, after all; the heart has its right, life its
poetry''. It is in vain to remind him that
personal worth is only imagination if there is nothing
eternal and unconditioned in the will: he dismisses it
as a subjective idea and seeks the riddle of life solved in
physiology. Unscientific, half-scientific
people take his unbelief for defiance. They do not understand
him; it is not defiance, it is conscience,
scientific conscience. It has perhaps cost
him sleepless nights and bitter tears before he let
go of the consoling illusions; but --- ``truth
is objective, science is objective!'' Would such a
graduate in modern knowledge feel a need to
let his culture go back ``into the dried-up ecclesiastical
channels'' and become a member of a neo-rationalist
congregation?

% --- p. 110 ---
\setcounter{footnote}{0}%
\opage{110}With the conscience of faith it is another matter;
it has a spark of eternity, something unconditioned in the self. But
precisely in this the opposition between the supernatural
and the natural comes forward in the innermost depths of consciousness.
When the personality is about to form itself, the spirit-will steps
up against the drive-will, and man finds himself
as a dualistic being. What is painful in
dualism does not consist in there being, in the ground of life
itself, dispositions of different kinds. The natural life of the individual furnishes
precisely the foundation of its spiritual life. No, the strife
arises from the drive-will, in its sensuous immediacy,
not being willing to submit to the spirit-will. What, then,
does the drive-will want? A momentary satisfaction of
natural desire or, if the momentary desire
is to be sacrificed, at least a satisfaction of a selfishly ruling
passion. And what does the spirit-will want? To rule,
and to rule unconditionally. It is a true torment for the
natural man to get the spirit-will along with him, for it
will never allow the drives, not even the gifted ones, to
break loose and go off on their own; it watches over
them like a guardian, it sees to their upbringing
and keeps them under strict discipline. Immediate
feeling may well find it outrageous
that the innocent inclinations, the merry genii of enjoyment,
these playing children of nature, are so pettily
kept in check; but conscience, the conscience of faith,
says the spirit-will is right; both agree that the unconditioned
must unconditionally be held fast. It is
not from lust for power, nor from hatred of the drives, nor
from contempt for the individual's natural dispositions, that the spirit-will
wants to rule in this way; no, it is for the sake of freedom,
for the sake of personality, so that there may come
% --- p. 111 ---
\setcounter{footnote}{0}%
\opage{111}independence into the individual's character, unity into man's
being. The true unity is unity of self: the goal of
development is the personal self.

By raising itself above the merely natural and transforming
the drives into means, the spirit-will shows its kinship
with the supernatural. The task of life cannot
be solved ascetically by subduing the natural; but neither can it
be solved monistically by denying the spirit-will
in order to cultivate the drive-will. If there is nothing at all
supernatural in man's will, then we must let
nature have its own way, and then things look doubtful
for the neo-rationalist principle of morality. What,
in the deepest sense, is the supernatural? Precisely
that in which everything natural is grounded. ``The more of the
supernatural there is in a human being's will, the more
personal is this will; the more personal the will is,
the more it will be able to work through the natural ground
of individuality, develop its faculties and strengthen
its dispositions, first and last its disposition to eternity, so that the
corruptible may die away, and the incorruptible
live.'' That is the testimony of the conscience of faith. ---
Would not a human being with such a ``faith of conscience''
prefer to remain in Luther's old Church
and --- with the reservation of the personal freedom of appropriation
--- bow in reverence before its dogmas and sacraments,
instead of seeking comfort and counsel in the rationalist
unity of faith and knowledge?

Existence is a wonder --- alles Sein ist Wunder;
it is admitted from the side of science as well as
from that of religion that here is a riddle that cannot
be solved. The present-day assumption that sensuous
nature with its phenomena and laws is a blind
% --- p. 112 ---
\setcounter{footnote}{0}%
\opage{112}product of a blind necessity that does not understand
itself is just as one-sided, and as unsatisfying for personal
self-understanding, as the Middle Ages'
notion of the arbitrary intervention of the supernatural
in the connection of things is outrageous to the present day's
scientific consciousness. The enigmatic interaction
of the supernatural and the natural --- the perennial
wonder --- is a unity-in-doubling of nature's means and
the spirit's ends, impenetrable to the human understanding, a unity-in-doubling which
from the standpoint of knowledge must be judged according to the laws of necessity,
from the standpoint of faith according to the laws of freedom.
As surely as science, on behalf of inquiry, criticism
and culture, has the right to declare hollow and fantastic every attempt, however
profound, that would explain the intervention of the supernatural
will by overlooking, misunderstanding or
subtly evading the laws of nature,
just as surely has religion, with its faith in revelation,
in the name of personality, humanity and freedom,
the right to protest against every encroachment that, under
the pretence of raising self-consciousness to a
higher level, makes the being of the self into semblance and knowingly
violates the laws of spirit. Within these limits
the measure is kept --- for a characterization of \emph{old}
and \emph{new prophets}.

\end{document}
